% concatenated from main.tex, setting.tex, abs.tex, intro.tex, sec2.tex, sec3.tex, sec4.tex, comps.tex, concl.tex, ref.tex
\documentclass{article}
\usepackage{setspace} % Allows the use of \toprule, \midrule and \bottomrule for better rules in tables
% \usepackage{biblatex}
% \addbibresource{references.bib}
% \usepackage{threeparttable}
% \usepackage{multirow}
% \usepackage{booktabs}       % For professional table formatting
% \usepackage{pdflscape}
\usepackage{booktabs}
\usepackage{longtable}
\usepackage{adjustbox}
\usepackage{geometry}
\usepackage{rotating} % For sideways tables
\geometry{a4paper, margin=1in}
\usepackage[english]{babel} % or polyglossia
\usepackage{csquotes}
\usepackage[backend=biber]{biblatex}
\addbibresource{references.bib}
\usepackage{amsthm} % Include this package
\usepackage{amsmath, amssymb} % Already included, needed for math symbols
\newcommand{\R}{\mathbb{R}} % Define \R as shorthand for \mathbb{R}
\newtheorem{corollary}{Corollary} 
\newcommand{\keywords}[1]{\textbf{Keywords: }#1}
\newtheorem{lemma}{Lemma}
\newtheorem{prop}{Proposition}
\theoremstyle{remark}
\newtheorem{remark}{Remark}
\newtheorem{theorem}{Theorem}
%\usepackage[english]{babel}
%\usepackage[linesnumbered]{algorithm2e}
\usepackage{algorithm}      % For algorithm environments
\usepackage{algpseudocode}  % For pseudocode within algorithms
\usepackage{graphicx} % Required for inserting images
\usepackage{xcolor}
\usepackage{graphicx}
\usepackage{amsmath, amsthm, amsfonts}
\usepackage{graphicx}
\usepackage{caption}
\usepackage[pdfborder={0 0 0},colorlinks=true,urlcolor=blue,citecolor=red,bookmarks=false]{hyperref}
%%%%%%%%%%%%%%
%  Cover Page
\newcommand{\tr}{\operatorname{tr}}
\usepackage{authblk}
\usepackage{siunitx}
\usepackage{booktabs}

\title{Two Generalized Derivative-free Methods to Solve Large Scale Nonlinear Equations with Convex Constraints}
\author[1,5]{Kabenge Hamiss}
\author[1,4]{Mohammed Alshahrani}
\author[2,3]{Mujahid N. Syed}
\affil[1]{\textit{Department of Mathematics, King Fahd University of Petroleum and
Minerals, Dhahran 31261, Saudi Arabia}}
\affil[2]{\textit{Department of Industrial \& Systems Engineering, King Fahd University of Petroleum and
Minerals, Dhahran 31261, Saudi Arabia}}
\affil[3]{\textit{Interdisciplinary Research Center for Intelligent Secure Systems, King Fahd University of Petroleum and
Minerals, Dhahran 31261, Saudi Arabia}}
\affil[4]{\textit{Interdisciplinary Research Center for Smart Mobility and Logistics, King Fahd University of Petroleum and
Minerals, Dhahran 31261, Saudi Arabia}}
\affil[5]{\textit{Department of Mathematics and Statistics, Islamic University in Uganda, Mbale 2555, Uganda}}



\begin{document}

\maketitle
%%%%%%%%%%%%%%%%%%%%%%%%
%%%%%%%%%%%%%%%%%%%%%%%%

% Abstract and Keywords
%\section*{Abstract}
\begin{abstract}%
    In this work, we propose two derivative-free methods to address the problem of large-scale nonlinear equations with convex constraints.  These algorithms satisfy the sufficient descent condition. The search directions can be considered generalizations of the Modified Optimal Perry conjugate gradient method and the conjugate gradient projection method or the Spectral Modified Optimal Perry conjugate gradient method and the Spectral Conjugate Gradient Projection method. The global convergence of the former does not depend on the Lipschitz continuity of G. In contrast, the latter's global convergence depends on the Lipschitz continuity of G. The numerical results show the efficiency of the algorithms.
\end{abstract}
\keywords{Spectral gradient, Descent methods, Line search methods, Derivative-free algorithms, Global Convergent algorithms}
%\keywords{Classical analysis, Acceleration, Convergence, Gradient descent, Algorithm, Unconstrained, Optimization} %Include a list of at least five keywords.
%\newpage
%\tableofcontents 

\vskip 0.5em
\hrulefill
%\newpage
\section{Introduction}
Let $\Gamma \subset \mathbb{R}^n$ be a nonempty, closed, and convex set and $G$ be continuous mapping such that,  
\begin{align}
    G(x) = 0. \label{eq:op}
\end{align}
where $x \in \Gamma$. Equations in \eqref{eq:op} are called a system of non-linear equations.\\
Nonlinear equations represent a significant class of problems closely related to optimization challenges, which frequently emerge in various science, technology, and industry fields. Many problems have been extensively studied, highlighting their importance in all disciplines \cite{ahookhosh2013two}. These problems are often complex, requiring sophisticated methods to solve efficiently.
The signals may be smooth or non-smooth and appear in a variety of applications, such as Bregman distances~\cite{iusem1997newton}, Monotone variational inequality problems ~\cite{he2002new, malitsky2014extragradient, zhao2001monotonicity}, Non-negative matrix factorization ~\cite{berry2007algorithms}, Phase retrieval~ \cite{candes2015phase}, Constrained neural networks~ \cite{chorowski2014learning}, Chemical equilibrium systems in thermodynamics ~ \cite{meintjes1987methodology, zeleznik1968calculation}, Financial modeling and forecasting ~ \cite{dai2020efficient}, Signal processing and signal restoration, especially in wavelet deconvolution and compressed sensing ~\cite{abubakar2020note, figueiredo2007gradient},   solid state physics, plasma and fluid mechanics \cite{chen2019gramian, gao2019mathematical}.\\ 
Given the wide range of applications, developing efficient algorithms to solve \eqref{eq:op} is crucial. Various methods have been proposed to solve these problems, with considerable research focusing on improving their accuracy and computational efficiency. Among such methods include methods involving determining the derivatives of $G(x)$ in \eqref{eq:op} such as Newton and Newton-like methods in ~\cite{iusem1997newton, nocedal1999numerical}. However, they require $G(x)$ to be not only continuous but also differentiable(smooth) which may not be the case in some instances. Further, even if $G$ is differentiable, it may be difficult to obtain derivatives of $G(x)$ in \eqref{eq:op} and this requires large memory for computation and storing. This has steered mathematicians to devise methods that do not require $G(x)$ to be differentiable. Such methods are called derivative-free methods. \\
Many derivative-free methods have been proposed stemming from Conjugate gradient methods(CGM) and Spectral conjugate gradient methods. They involve determining a descent direction for subsequent iterations to be close to the solution. The appropriate step length accompanies this. Constructing a hyperplane for the current iterate to be separated from the solution set due to convex separation. Projecting $x_k$ onto the hyperplane for the algorithm's convergence.\\
Sabi'u et al.~\cite{sabi2023modified} proposed a modified optimal Perry conjugate gradient. \\method(MOPCGM) with conjugate parameter
\begin{equation}
   \theta_k^{MOP}=\frac{(v_{k-1}-\zeta_k^*s_{k-1})^TG_k}{p_{k-1}^Tv_{k-1}} 
\end{equation}
 where $\zeta_k^*=\frac{s_{k-1}^Tw_{k-1}}{||s_{k-1}||^2}$,
 \begin{equation}
     z_{k-1}=x_{k-1}+\alpha_{k-1}p_{k-1},
 \end{equation}
\begin{equation}
    s_{k-1}=z_{k-1}-x_{k-1}
\end{equation}
and
\begin{equation}
    v_{k-1}=G_k-G_{k-1}+\tau s_{k-1}, ~~~\tau\geq 0
\end{equation}

In 2020, Zheng et al~\cite{zheng2020conjugate} proposed a Conjugate gradient projection method (CGPM) with conjugate parameter 
\begin{equation}
\label{eq:HZ}
\theta_k=\frac{G_k^Tp_{k-1}-2a_k||w_{k-1}||^2}{p_{k-1}^Tw_{k-1}}
\end{equation}
  $a_k=\frac{G_k^Tp_{k-1}}{w_{k-1}^Tp_{k-1}}$, $w_{k-1}=y_{k-1}+t_{k-1}p_{k-1}$, $t_{k-1}=\max\{1,1-\frac{p_{k-1}^Ty_{k-1}}{||p_{k-1}||^2}\}$, and  $\gamma, \tau>0$\\
 % In this work, we propose two methods that generalize the Modified Optimal Perry Conjugate Gradient Method (MOPCGM) and the Conjugate Gradient Projection Method (CGPM). The main contributions of this paper include;
 % \begin{enumerate}
 %     \item[(i)] The new methods outperform the former in terms of the number of function evaluations
 %     \item[(ii)]  They take fewer iterations than the former
 %     \item[(iii)] They also require less time to solve the problem compared to the former. 
 %     \item[(iv)] They are also more accurate and robust compared to the former. 
 % \end{enumerate}
In this work, we propose two novel methods that generalize the Modified Optimal Perry Conjugate Gradient Method (MOPCGM) and the Conjugate Gradient Projection Method (CGPM). The main contributions of this paper are summarized as follows: \begin{enumerate} \item[(i)] The proposed methods outperform the original ones in terms of the number of function evaluations. \item[(ii)] They converge in fewer iterations. \item[(iii)] They require less computational time to solve the problem. \item[(iv)] They exhibit improved accuracy and robustness. \end{enumerate}
The paper is arranged as follows: In the next section, we derive the generalized MOPCGM and CGPM. Section 3 details the numerical experiments and their application to compressed sensing. In the last section, we give the conclusion.
\section{Generalized Derivative-Free Methods}
\label{chp:generalMOPCGM}

{In this section, we give a generalized result for the Perry conjugate gradient and Conjugate gradient projection techniques. \\
We also detail the convergence analysis of algorithms  provided the following assumptions hold:
\begin{enumerate}
    \item[(A1)] $\Gamma$ is nonempty
    \item[(A2)] G is Lipschitz continuous. Then there exists some nonzero constant $L$ such that for any $x, y\in \mathbf{R}^n$ then $||G(y)-G(x)||\leq L||y-x||$
\end{enumerate}
\subsection{Generization of MOPCGM}
From $s_{k-1}=x_k-x_{k-1}$, $g_k=\nabla G(x_k)$ and $y_{k-1}=g_k-g_{k-1}$. Also assume that $B_{k-1}$ is positive-definite, then 
\[
B_k=B_{k-1}-\frac{B_{k-1}s_{k-1}(B_{k-1}s_{k-1})^T}{s_{k-1}^TB_{k-1}s_{k-1}}+\frac{y_{k-1}y_{k-1}^T}{s_{k-1}^Ty_{k-1}}
\]
It is a rank 2 matrix update called DFP.

Now, using the Quasi-Newton  approach as used in ~\cite{ perry1978modified, sabi2023modified} we propose 
\begin{equation}
  p_k=-\Tilde{Q_k}G_k  
\end{equation}
such that
\begin{equation}\label{eq:update}
  \Tilde{Q_k}=\lambda I-\lambda\frac{y_{k-1}s_{k-1}^T}{2y_{k-1}^Ts_{k-1}}-\lambda\frac{s_{k-1}y_{k-1}^T}{2y_{k-1}^Ts_{k-1}}+t_k\frac{s_{k-1}s_{k-1}^T}{y_{k-1}^Ts_{k-1}}  
\end{equation}
$\lambda>0$ and $t_k>0$. Because $G$ is monotone, then $s_{k-1}^Ty_{k-1}> 0$ for all $x \neq x^*$. This implies that both $s_{k-1}$ and $y_{k-1}$ are nonzero vectors. Let $D$ be spanned by $\{s_{k-1}, y_{k-1} \}$ and $a$ be any vector in $\R^n$ such that $a^TD\neq 0$, then 
\[
a^T\Tilde{Q_k}a=t_k\frac{(a^Ts_{k-1})^2}{y_{k-1}^Ts_{k-1}}>0
\]
and this shows that $\Tilde{Q_k}$ is positive definite.\\
But $\Tilde{Q_k}$ is rank 2 matrix update, then eigenvalue $\lambda$ is of multiplicity $n-2$.
Now we need to determine the other two eigenvalues $\eta_k^-$ and $\eta_k^+$ since $\Tilde{Q_k}$ is full rank based on the symmetric property of $\Tilde{Q_k}$. Therefore we can find a set of vectors that are mutually orthogonal $\{u_k^i\}_{i=1}^{n-2}$ such that \[
\Tilde{Q}_ku_k^i=\lambda u_k^i, ~~~i=1, \cdots, n-2
\] 
and satisfy ${u_k^i}^TD=0$ for $i=1, \cdots, n-2$.\\
Therefore $\{u_k^i\}_{i=1}^{n-2}$ are eigenvectors of $\Tilde{Q}_k$ with the corresponding eigenvalue $\lambda$ for every $u_k^i$. We can now determine the remaining $2$ eigenvalues of $\Tilde{Q}_k$ that is $\eta_k^+$ and $\eta_k^-$. The following lemmas are crucial.
% Because the trace of a square matrix is the sum of its eigenvalues, we have the following Lemmas.
\begin{lemma}
\label{A1}
    Let $\Tilde{Q_k}$ be defined as in \eqref{eq:update}, then \[
    \tr{\Tilde{Q_k}}=\lambda (n-1)+t_k\frac{||s_{k-1}||^2}{s_{k-1}^Ty_{k-1}}\] 
\end{lemma}
\begin{proof}
    Using the linearity property of the trace of a matrix, it is the sum of all eigenvalues of a matrix. So we have 
    \[
    \tr{\Tilde{Q_k}}=\tr{\lambda I}-\tr{\lambda\frac{y_{k-1}s_{k-1}^T}{2y_{k-1}^Ts_{k-1}}}-\tr{\lambda\frac{s_{k-1}y_{k-1}^T}{2y_{k-1}^Ts_{k-1}}}+\tr{t_k\frac{s_{k-1}s_{k-1}^T}{y_{k-1}^Ts_{k-1}}}
    \]
    but $\tr{\lambda I}=n\lambda$. Let $A=y_{k-1}s_{k-1}^T$, $A^T=s_{k-1}y_{k-1}^T$ and $B=s_{k-1}s_{k-1}^T$  then $\tr{A}=\tr{A^T}=y_{k-1}^Ts_{k-1}$ and $\tr{B}=||s_{k-1}||^2$. Then 
    \[
    \tr{\Tilde{Q_k}}=\lambda(n-1)+t_k\frac{||s_{k-1}||^2}{y_{k-1}^Ts_{k-1}}
    \]
\end{proof}

\begin{lemma}
\label{A2}
   Let $\Tilde{Q_k}$ be defined as in \eqref{eq:update}, then
  \[
  \tr{\Tilde{Q}_k^T\Tilde{Q}_k}=\lambda^2( n-\frac{3}{2})+\frac{\lambda^2}{2}\frac{||s_{k-1}||^2||y_{k-1}||^2}{(s_{k-1}^Ty_{k-1})^2}+t_k^2\frac{||s_{k-1}||^4}{(s_{k-1}^Ty_{k-1})^2}\]
\end{lemma}
\begin{proof}
    From 
    \begin{multline*}
    \Tilde{Q}_k^T\Tilde{Q}_k=\left(\lambda I-\lambda\frac{y_{k-1}s_{k-1}^T}{2y_{k-1}^Ts_{k-1}}-\lambda\frac{s_{k-1}y_{k-1}^T}{2y_{k-1}^Ts_{k-1}}+t_k\frac{s_{k-1}s_{k-1}^T}{y_{k-1}^Ts_{k-1}}\right)^T\\
    \left(\lambda I
    -\lambda\frac{y_{k-1}s_{k-1}^T}{2y_{k-1}^Ts_{k-1}}-\lambda\frac{s_{k-1}y_{k-1}^T}{2y_{k-1}^Ts_{k-1}}+t_k\frac{s_{k-1}s_{k-1}^T}{y_{k-1}^Ts_{k-1}}\right)
    \end{multline*}
    
   \begin{multline*}
    \Tilde{Q}_k^T\Tilde{Q}_k=\left(\lambda I-\lambda\frac{s_{k-1}y_{k-1}^T}{2y_{k-1}^Ts_{k-1}}-\lambda\frac{y_{k-1}s_{k-1}^T}{2y_{k-1}^Ts_{k-1}}+t_k\frac{s_{k-1}s_{k-1}^T}{y_{k-1}^Ts_{k-1}}\right)\\
    \left(\lambda I
    -\lambda\frac{y_{k-1}s_{k-1}^T}{2y_{k-1}^Ts_{k-1}}-\lambda\frac{s_{k-1}y_{k-1}^T}{2y_{k-1}^Ts_{k-1}}+t_k\frac{s_{k-1}s_{k-1}^T}{y_{k-1}^Ts_{k-1}}\right)
    \end{multline*}

   \begin{multline*}
    \Tilde{Q}_k^T\Tilde{Q}_k=\lambda^2 I-\lambda^2\frac{3s_{k-1}y_{k-1}^T}{4y_{k-1}^Ts_{k-1}}-\lambda^2\frac{3y_{k-1}s_{k-1}^T}{4y_{k-1}^Ts_{k-1}}+t_k\lambda\frac{s_{k-1}s_{k-1}^T}{y_{k-1}^Ts_{k-1}}+\lambda^2\frac{||s_{k-1}||^2y_{k-1}y_{k-1}^T}{4(y_{k-1}^Ts_{k-1})^2}\\-t_k\lambda\frac{||s_{k-1}||^2y_{k-1}s_{k-1}^T}{2(y_{k-1}^Ts_{k-1})^2}+\lambda^2\frac{||y_{k-1}||^2s_{k-1}s_{k-1}^T}{4(y_{k-1}^Ts_{k-1})^2}-\lambda t_k\frac{||s_{k-1}||^2s_{k-1}y_{k-1}^T}{2(y_{k-1}^Ts_{k-1})^2}+t_k^2\frac{||s_{k-1}||^2s_{k-1}s_{k-1}^T}{(y_{k-1}^Ts_{k-1})^2}
    \end{multline*}
    we have seen that $\tr{y_{k-1}s_{k-1}^T}=y_{k-1}^Ts_{k-1}$ and $\tr{s_{k-1}s_{k-1}^T}=||s_{k-1}||^2$. Therefore
    properties of trace of matrices
    \begin{multline*}
    \tr{\Tilde{Q}_k^T\Tilde{Q}_k}=\tr{\lambda^2 I}-\tr{\lambda^2\frac{3s_{k-1}y_{k-1}^T}{4y_{k-1}^Ts_{k-1}}}-\tr{\lambda^2\frac{3y_{k-1}s_{k-1}^T}{4y_{k-1}^Ts_{k-1}}}+\tr{t_k\lambda\frac{s_{k-1}s_{k-1}^T}{y_{k-1}^Ts_{k-1}}}+\tr{\lambda^2\frac{||s_{k-1}||^2y_{k-1}y_{k-1}^T}{4(y_{k-1}^Ts_{k-1})^2}}\\-\tr{t_k\lambda\frac{||s_{k-1}||^2y_{k-1}s_{k-1}^T}{2(y_{k-1}^Ts_{k-1})^2}}+\tr{\lambda^2\frac{||y_{k-1}||^2s_{k-1}s_{k-1}^T}{4(y_{k-1}^Ts_{k-1})^2}}-\tr{\lambda t_k\frac{||s_{k-1}||^2s_{k-1}y_{k-1}^T}{2(y_{k-1}^Ts_{k-1})^2}}+\tr{t_k^2\frac{||s_{k-1}||^2s_{k-1}s_{k-1}^T}{(y_{k-1}^Ts_{k-1})^2}}
   \end{multline*}
   
Therefore

   \begin{multline*}
   \tr{\Tilde{Q}_k^T\Tilde{Q}_k}=n\lambda^2 -\lambda^2\frac{3}{4}-\lambda^2\frac{3}{4}+t_k\lambda\frac{||s_{k-1}||^2}{y_{k-1}^Ts_{k-1}}+\lambda^2\frac{||s_{k-1}||^2||y_{k-1}||^2}{4(y_{k-1}^Ts_{k-1})^2}\\-t_k\lambda\frac{||s_{k-1}||^2}{2(y_{k-1}^Ts_{k-1})}+\lambda^2\frac{||y_{k-1}||^2||s_{k-1}||^2}{4(y_{k-1}^Ts_{k-1})^2}-\lambda t_k\frac{||s_{k-1}||^2}{2(y_{k-1}^Ts_{k-1})}+t_k^2\frac{||s_{k-1}||^4}{(y_{k-1}^Ts_{k-1})^2}
   \end{multline*}

 \[
   \tr{\Tilde{Q}_k^T\Tilde{Q}_k}=\lambda^2(n-\frac{3}{2})+\lambda^2\frac{||s_{k-1}||^2||y_{k-1}||^2}{2(y_{k-1}^Ts_{k-1})^2}+t_k^2\frac{||s_{k-1}||^4}{(y_{k-1}^Ts_{k-1})^2}
   \]
   for all $\lambda$.
\end{proof}
\begin{lemma}
    \label{lemma:prodegnval}
    The product of eigenvalues $\eta^+$ and $\eta^-$ is given by $\eta^+\eta^-=\frac{\lambda^2}{4}-\frac{\lambda^2}{4}\left(\frac{||s_{k-1}||||y_{k-1}||}{s_{k-1}^Ty_{k-1}}\right)^2+\lambda t_k\frac{||s_{k-1}||^2}{s_{k-1}^Ty_{k-1}}$.
\end{lemma}
\begin{proof}
Using Lemma \ref{A1} and the usual sum of all eigenvalues as the trace of a matrix, we get
\[
\lambda (n-2)+\eta_k^-+\eta_k^+=\lambda (n-1)+t_k\frac{||s_{k-1}||^2}{s_{k-1}^Ty_{k-1}}.
\]
Also using Lemma \ref{A2}, we obtain
\[
\lambda^2 (n-2)+{\eta_k^-}^2+{\eta_k^+}^2=\lambda^2( n-\frac{3}{2})+\frac{\lambda^2}{2}\frac{||s_{k-1}||^2||y_{k-1}||^2}{(s_{k-1}^Ty_{k-1})^2}+t_k^2\frac{||s_{k-1}||^4}{(s_{k-1}^Ty_{k-1})^2}.
\]
 Implying that
 \begin{equation}
     {\eta_k^-}^2+{\eta_k^+}^2=\frac{\lambda^2}{2}+\frac{\lambda^2}{2}\frac{||s_{k-1}||^2||y_{k-1}||^2}{(s_{k-1}^Ty_{k-1})^2}+t_k^2\frac{||s_{k-1}||^4}{(s_{k-1}^Ty_{k-1})^2}.
 \end{equation}
 Let $a=\frac{||s_{k-1}||^2}{s_{k-1}^Ty_{k-1}}$ and $b=\frac{||s_{k-1}||||y_{k-1}||}{s_{k-1}^Ty_{k-1}}$, then we obtain
 \begin{equation}\label{gmop3}
     \eta_k^-+\eta_k^+=\lambda+at_k
 \end{equation} and 
 \begin{equation}\label{gmop4}
     {\eta_k^-}^2+{\eta_k^+}^2=\frac{\lambda^2}{2}+\frac{\lambda^2}{2}b^2+a^2t_k^2
 \end{equation}respectively.\\
 Now from \eqref{gmop3} and \eqref{gmop4}, we obtain
 \begin{equation}\label{gmop5}
     \eta_k^-\eta_k^+=\frac{\lambda^2}{4}-\frac{\lambda^2}{4}b^2+\lambda at_k.
 \end{equation}
 \begin{remark}
     when $\lambda=1$, then we obtain 
 \[
  \eta_k^-\eta_k^+=\frac{1}{4}-\frac{1}{4}b^2+at_k
 \]
 which is the case in {MOPCGM} ~\cite{sabi2023modified}.
 \end{remark}
\end{proof}


\begin{lemma}
\label{lemma:optimalt}    
Let $\Tilde{Q}_k$ be defined as in \eqref{eq:update}, then the optimal value of $t$ is $t_k^*=\lambda\frac{s_{k-1}^Ty_{k-1}}{||s_{k-1}||^2}$.
\end{lemma}

 
 From \eqref{gmop3} and \eqref{gmop5}, we obtain
 \begin{equation}
     \eta_k^2-(\lambda+at)\eta_k+\left[\frac{\lambda^2}{4}-\frac{\lambda^2}{4}b^2+at_k\lambda\right]=0.
 \end{equation}
 This means
 
 \begin{equation}
     \eta_k^{\pm}=\frac{(\lambda+at_k)\pm\sqrt{(\lambda+at)^2 - 4(\frac{\lambda^2}{4}-\frac{\lambda^2}{4}b^2+at_k\lambda)}}{2}.
 \end{equation}
 For positive definiteness of the matrix $\Tilde{Q}_k$, by applying some algebra, $t>\frac{\lambda}{4a}(b^2-1)$.\\
 Consequently, we obtain 
 \[
 t>\frac{\lambda}{4}(\frac{||y_{k-1}||^2}{s_{k-1}^Ty_{k-1}}-\frac{s_{k-1}^Ty_{k-1}}{||s_{k-1}||^2}).
 \]
 Now to obtain the optimal value of $t_k$, we minimize the square of the difference between $\eta_k^+$ and $\eta_k^-$ so that the condition number is small, then 
 \begin{equation}
   (\eta_k^+-\eta_k^-)^2 =(\lambda+at)^2 - 4(\frac{\lambda^2}{4}-\frac{\lambda^2}{4}b^2+at_k\lambda).
 \end{equation}
Therefore, the optimal value of $t$ is $\frac{\lambda }{a}$ .\\
That is 
\[
t_k^*=\lambda\frac{s_{k-1}^Ty_{k-1}}{||s_{k-1}||^2}.
\]
\begin{remark}
For $\lambda=1$, we obtain the Optimal Perry parameter $t_k^*=\frac{s_{k-1}^Ty_{k-1}}{||s_{k-1}||^2}$ obtained by Subu'i et al~\cite{sabi2023modified}.
\end{remark}

 We propose a new search direction is given by 
\begin{equation}
p_k=\begin{cases}\label{GMOP}
    -G_k & k=0\\
    -M_kG_k+\theta_k^{GMOP}p_{k-1} & k\geq 1
\end{cases}
\end{equation}
such that $M_k=\lambda +\theta_k^{GMOP}\frac{G_k^Tp_{k-1}}{||G_k||^2}$ and $\theta_k^{GMOP}=\frac{(v_{k-1}-t_k^*s_{k-1})^TG_k}{p_{k-1}^Tv_{k-1}}$ and $t_k^*=\lambda \frac{s_{k-1}^Tv_{k-1}}{||s_{k-1}||^2}$, $\lambda > 0$, $v_k=y_{k-1}+\tau s_{k-1}$ and $\tau>0$. 
\begin{remark}
    \label{rem:adaptive}
    To adaptively choose $\lambda$ depending on the problem and the current iterate, we define $\lambda_k$ as the projection of two quantities onto the interval $[\alpha_{\min}, \alpha_{\max}]$, where $\alpha_{\min} > 0$. This ensures that $\lambda_k$ remains bounded between $\alpha_{\min}$ and $\alpha_{\max}$, adjusting itself dynamically based on the progress of the optimization.

We define:
\[
    \lambda_k = \Pi_{[\alpha_{\min}, \alpha_{\max}]}\left( \frac{\|v_{k-1}\|^2}{s_{k-1}^\top v_{k-1}}, \frac{s_{k-1}^\top v_{k-1}}{\|s_{k-1}\|^2} \right),
\]

The projection operator $\Pi_{[\alpha_{\min}, \alpha_{\max}]}(\cdot)$ projects the values onto the interval $[\alpha_{\min}, \alpha_{\max}]$.\\
   $\frac{\|v_{k-1}\|^2}{s_{k-1}^\top v_{k-1}}$ measures how much the $G(x)$ changes relative to the step size. A large value indicates a significant change in gradients, suggesting that a larger $\lambda_k$ may be beneficial.\\
   $ \frac{s_{k-1}^\top v_{k-1}}{\|s_{k-1}\|^2}$
   measures the alignment between the step and the change in $G(x)$, relative to the step length. A smaller value suggests a smaller $\lambda_k$.

%### Projection onto the Interval:
\begin{equation}
    \label{eq:adaptivegmopcgm}
\lambda_k = \Pi_{[\alpha_{\min}, \alpha_{\max}]}\left(\max\left(\frac{\|v_{k-1}\|^2}{s_{k-1}^\top v_{k-1}}, \frac{s_{k-1}^\top v_{k-1}}{\|s_{k-1}\|^2}\right)\right)
\end{equation}
This ensures that $\lambda_k$ is bounded within $[\alpha_{\min}, \alpha_{\max}]$, with $\alpha_{\min} > 0$ to prevent the step size from becoming too small.

%### Choosing $\alpha_{\min}$ and $\alpha_{\max}$:
- $\alpha_{\min}$ ensures that $\lambda_k$ does not become too small, which could cause slow convergence.
- $\alpha_{\max}$ limits $\lambda_k$ from growing too large, preventing instability.

%### Conclusion:
This adaptive scheme allows $\lambda_k$ to tune itself according to the characteristics of the problem, without requiring manual adjustments for each new optimization problem. The dynamic nature of $\lambda_k$ ensures a balance between stability and fast progress.

\end{remark}
Using \eqref{eq:adaptivegmopcgm}, \eqref{GMOP} becomes
\begin{equation}
p_k=\begin{cases}\label{eq:adaptGMOP}
    -G_k & k=0\\
    -M_kG_k+\theta_k^{GMOP}p_{k-1} & k\geq 1
\end{cases}
\end{equation}
such that $M_k=\lambda_k +\theta_k^{GMOP}\frac{G_k^Tp_{k-1}}{||G_k||^2}$ and 
\begin{equation}
    \label{eq:optimalconjpamop}
    \theta_k^{GMOP}=\frac{(v_{k-1}-t_k^*s_{k-1})^TG_k}{p_{k-1}^Tv_{k-1}},
\end{equation}

\begin{equation}
    \label{eq:optimaltmop}
    t_k^*=\lambda_k \frac{s_{k-1}^Tv_{k-1}}{||s_{k-1}||^2},
\end{equation} and
 $\lambda_k > 0$ defined in \eqref{rem:adaptive}, $v_k=y_{k-1}+\tau s_{k-1}$ and $\tau>0$. In lemma \ref{lemm1:gmopcgm}, we verified that the $p_k$ of \eqref{GMOP} satisfies the sufficient descent condition.
\begin{lemma}
\label{lemm1:gmopcgm}
Let $\{p_k\}$ and $G_k$ be produced by algorithm \ref{alg:gmopcg}. Let $\alpha_{min} > 0$, then  $G_k^Tp_k\leq-\alpha_{min}  ||G_k||^2.$ 
\end{lemma}
\begin{proof}
    To verify this, we multiply $G_k^T$ in \eqref{GMOP} and we obtain
\[
G_k^Tp_k=-\lambda_k ||G_k||^2-\theta_k^{GMOP}\frac{G_k^Tp_{k-1}}{||G_k||^2}||G_k||^2+\theta_k^{GMOP}G_k^Tp_{k-1},
\]
\[
G_k^Tp_k=-\lambda_k ||G_k||^2, 
\]
but from \eqref{eq:adaptivegmopcgm} $\lambda_k\geq \alpha_{min}$. It implies that 
\begin{equation}
\label{eq:sufficientdescentmop}
G_k^Tp_k\leq -\alpha_{min} ||G_k||^2, 
\end{equation}
\end{proof}



From the Lemma \ref{lemm1:gmopcgm} above, the descent condition generally does not depend on $t_k$.


%%%%%%%%%%%%%%%%%%%%%%%%%%%%%%%%%%%%%%%%%%%%%%%%%%%%%%%%%%%%%%%%%%%%%%%%%%%%%%%%%%%%%%%%%%%%%%%%%%%%%%%%%%%

%% Please enter the Section Title in the following in place of xxxxx. Use a label in place of yyyyy to be
%% referred later.
%% Please write the text in between the braces in place of AAAAA
\newpage
\subsubsection{Generalized MOPCGM Algorithm}
\renewcommand{\baselinestretch}{0.95}
\begin{algorithm}[H]
    \caption{Generalized MOPCGM}
    \label{alg:gmopcg}
    \begin{algorithmic}[1]
        \State \textbf{Input:} Function $G$, initial guess $x_0 \in \mathbb{R}^n$, 
        $\lambda_o > 0$, tolerance $\epsilon = \text{Tol}$, 
        $\rho \in (0, 1)$, parameters $\tau$,   
        $\beta > 0$, $\eta > 0$, $\zeta > 0$, $\alpha_{min}> 0$, $\alpha_{max}> \alpha_{min}$, $\gamma, \gamma_1, \gamma_2, \gamma_3, \gamma_4 \in (0, 2) $, $0<\zeta_1\leq \zeta_2$, projection function 
        $\Pi_{\Gamma}$ on convex set $\Gamma$
        \State \textbf{Output:} Solution $x^*$
        \State \textbf{Initialization:} Set $k \gets 0$, $\lambda \gets \lambda_o$, $z_k \gets x_0$, initialize $\alpha_k$
        
        \While{$\| G_k \| > \epsilon$}
             \If{$\|G(x_{k+1})\| < \|G(x_k)\|$ }
                \State $\lambda_{k+1} \gets \lambda_k$
                \State \textbf{Else}
                {$\lambda_{k+1} \gets \Pi_{[\alpha_{min}, \alpha_{max}]}(\max(\frac{\|v_{k-1}\|^2}{s_{k-1}^Tv_{k-1}}, \frac{s_{k-1}^Tv_{k-1}}{\|s_{k-1}\|^2}))$}
            \EndIf
            \State Determine $p_k$ as in \eqref{GMOP}
            \State Adjust $\alpha_k = \max \{ \rho^i \beta \}$ such that 
            \begin{equation}
                G(x_k + \alpha_k p_k)^T p_k \leq -\zeta \alpha_k \| p_k \|^2 \Pi_{[\zeta_1, \zeta_2]}(\|G(x_k + \alpha_k p_k)\|)
                \label{eq:backtrack}
            \end{equation}
            \If{$z_k = x_k + \alpha_k p_k \in \Gamma$ and $\| G(z_k) \| < \epsilon$}
                \State $x^* \gets z_k$
                \State \textbf{break}
            \EndIf
            
            \State Compute $\mu_k \gets \frac{G(z_k)^T (x_k - z_k)}{\| G(z_k) \|^2}$
            \State Update $x_{k+1} \gets \Pi_{\Gamma} (x_k - \gamma\mu_k G(z_k))$
            \State Compute $s_{k-1}=z_{k-1}-x_{k-1} $ and $v_{k-1}=G(x_{k})-G(x_{k-1})+\tau s_{k-1}$
            \State Compute $\theta_k^{MOP} $ and $t^*$ using \eqref{eq:optimaltmop}  and \eqref{eq:optimalconjpamop} respectively. 
            \If{$\|f_{k+1}\| < \|f_k\|$}
    \State $\gamma = \min(\gamma \cdot \gamma_1, \gamma_2)$
\Else
    \State $\gamma = \max(\gamma \cdot \gamma_3, \gamma_4)$
    \State \textbf{break}
\EndIf

            \If{$\|p_k\|  \approx 0$ }
                \State $x^* \gets x_k$
                \State \textbf{break}
            \EndIf   
            \State Set $x_k \gets x_{k+1}$
        \EndWhile
        \State \textbf{Return} $x^*$
    \end{algorithmic}
\end{algorithm}
\renewcommand{\baselinestretch}{1.50}

% On this basis, we will point out the relationship between MOPCGM and CGPM.
}
%%%%%%%%%%%%%%%%%%%%%%%%%%%%%%%%%%%%%%%%%%%%%%%%%%%%%%%%%%%%%%%%%%%%%%%%%%%%%%%%%%%%%%%%%%%%%%%%%%%%%%%%%%%

%% Please enter the Section Title in the following in place of xxxxx. Use a label in place of yyyyy to be
%% referred later.
%% Please write the text in between the braces in place of AAAAA
\newpage

\subsubsection{Convergence Analysis of Generalized MOPCGM}
\label{GeneralizedConvergence}

{The next Lemma examines that the algorithm is well defined.
\begin{lemma}
\label{lemm2:gmopcgm}
    Let $G$ be Lipschitz continuous, then 
    \[
    \alpha_k\geq \min\{\eta, \frac{\alpha_{min}\rho}{L+\zeta\zeta_2 }\frac{||G_k||^2}{||p_k||^2}\}.
    \]
\end{lemma}
\begin{proof}
    With the line search in  Algorithm \ref{alg:gmopcg}, Let $\alpha_k$ be the optimal step length that satisfies \eqref{eq:backtrack}, then $\Tilde{\alpha}_k = \alpha_k/\rho $ violets \eqref{eq:backtrack}. So, 
\[
G(x_k+\tilde{\alpha}_kp_k)^Tp_k>- \zeta \tilde{\alpha}_k||p_k||^2\Pi_{[\zeta_1, \zeta_2]}(\|G(x_k + \Tilde{\alpha}_k p_k)\|)
\]
and  $G_k^Tp_k\leq -\alpha_{min}||G_k||^2$. 
\[
||G_k||^2\leq \frac{1}{\alpha_{min}}\left[(G(x_k+\Tilde{\alpha}_kp_k)- G_k)^Tp_k-G(x_k+\Tilde{\alpha}_kp_k)^Tp_k\right]
\]
\[
||G_k||^2\leq \frac{1}{\alpha_{min}}[(G(x_k+\Tilde{\alpha}_kp_k)- G_k)^Tp_k+\zeta \tilde{\alpha}_k||p_k||^2\zeta_2]
\]

\[
||G_k||^2\leq \frac{\alpha_k(L+\zeta_2\zeta) }{\rho\alpha_{min}}||p_k||^2
\]
This completes the proof.
\end{proof}

\begin{lemma}
\label{lemm4:gmopcgm}
     The direction $p_k$ produced by  Algorithm \ref{alg:gmopcg} meets the trust region property
    \[
    \alpha_{min} \|G(x_k)\|\leq \|p_k\|\leq \kappa \|G(x_k)\|.
    \]
\end{lemma}
\begin{proof}

From 
\[
s_{k-1}=z_{k-1}-x_{k-1}=\alpha_{k-1}p_{k-1}
\]
Also $|v_{k-1}^Ts_{k-1}|\leq (L\gamma+\tau)||s_{k-1}||^2 $,  then $|t_k^{GMOP}|= \lambda_k \frac{|v_{k-1}^Ts_{k-1}|}{||s_{k-1}||^2}\leq \alpha_{max}(L\gamma+\tau)$.

 Therefore
\[
|\theta_k^{GMOP}|=|\frac{v_{k-1}^TG_k}{p_{k-1}^Tv_{k-1}}-\lambda_k \frac{s_{k-1}^Tv_{k-1}}{||s_{k-1}||^2}\frac{s_{k-1}^TG_k}{p_{k-1}^Tv_{k-1}}|,
\]
\[
|\theta_k^{GMOP}|\leq \frac{|v_{k-1}^TG_k|}{\tau\alpha_{k-1}||p_{k-1}||^2}+ \frac{\alpha_{max}(L\gamma+\tau)||s_{k-1}||^2}{||s_{k-1}||^2}\frac{|s_{k-1}^TG_k|}{\tau\alpha_{k-1}||p_{k-1}||^2},
\]
\[
|\theta_k^{GMOP}|\leq \frac{||v_{k-1}||||G_k||}{\tau\alpha_{k-1}||p_{k-1}||^2}+ \frac{\alpha_{max}(L\gamma+\tau)||s_{k-1}||^2}{||s_{k-1}||^2}\frac{||s_{k-1}||||G_k||}{\tau\alpha_{k-1}||p_{k-1}||^2},
\]
\[
|\theta_k^{GMOP}|\leq \frac{(L\gamma+\tau)\alpha_{k-1}||p_{k-1}||||G_k||}{\tau\alpha_{k-1}||p_{k-1}||^2}+\alpha_{max} \frac{(L\gamma+\tau)||s_{k-1}||^2}{||s_{k-1}||^2}\frac{\alpha_{k-1}||p_{k-1}||||G_k||}{\tau\alpha_{k-1}||p_{k-1}||^2},
\]
\[
|\theta_k^{GMOP}|\leq \frac{(L\gamma+\tau)||G_k||}{\tau||p_{k-1}||}+\alpha_{max}(L\gamma+\tau) \frac{||G_k||}{\tau||p_{k-1}||},
\]
\begin{equation}\label{mest1}
 |\theta_k^{GMOP}|\leq (1+\alpha_{max})\frac{(L\gamma+\tau)||G_k||}{\tau||p_{k-1}||} ,  
\end{equation}
 Now 
  \[ ||p_k||\leq |\lambda_k-\theta_k^{GMOP}\frac{G_k^Tp_{k-1}}{||G_k||^2}|||G_k||+|\theta_k^{GMOP}|||p_{k-1}||, \]
 \[
 ||p_k||\leq |\lambda-\theta_k^{GMOP}|||G_k||+|\theta_k^{GMOP}\frac{G_k^Tp_{k-1}}{||G_k||^2}|||p_{k-1}||.
 \]
 But also 
 \[
 |(\lambda_k-\theta_k^{GMOP}\frac{G_k^Tp_{k-1}}{||G_k||^2})G_k+\theta_k^{GMOP}p_{k-1}|\leq \alpha_{max}||G_k||+|\theta_k^{GMOP}|\frac{G_k^Tp_{k-1}}{||G_k||^2}|||G_k||+|\theta_k^{GMOP}|||p_{k-1}||.
 \]
\[
 |(\lambda_k-\theta_k^{GMOP}\frac{G_k^Tp_{k-1}}{||G_k||^2})G_k+\theta_k^{GMOP}p_{k-1}|\leq \alpha_{max}||G_k||+|\theta_k^{GMOP}|\frac{\|G_k\|\|p_{k-1}\|}{||G_k||^2}|||G_k||+|\theta_k^{GMOP}|||p_{k-1}||.
 \]
 It implies that
\[
 \|p_k\|\leq \alpha_{max}||G_k||+|\theta_k^{GMOP}|\|p_{k-1}\|+|\theta_k^{GMOP}|||p_{k-1}||,
 \]
Therefore,
\[
 \|p_k\|\leq \alpha_{max}||G_k||+2|\theta_k^{GMOP}|\|p_{k-1}\|.
 \]
 Using \eqref{mest1}
 \[
 \|p_k\|\leq \alpha_{max}||G_k||+2(1+\alpha_{max})\frac{(L\gamma+\tau)||G_k||}{\tau||p_{k-1}||} \|p_{k-1}\|.
 \]
\[
 \|p_k\|\leq \alpha_{max}||G_k||+2(1+\alpha_{max})\frac{(L\gamma+\tau)||G_k||}{\tau} .
 \]
We finally get 
\begin{equation}
\label{eq:upperboundgmopcgm}
\|p_k\| \leq \left(\alpha_{max}+2(1+\alpha_{max})\frac{(L\gamma+\tau)}{\tau}\right)\|G_k\| .
\end{equation}
From \eqref{eq:upperboundgmopcgm}, we conclude that
\begin{equation}\label{mest4}
||p_k||\leq \kappa ||G_k||,
\end{equation}
where 
\begin{equation}
\label{eq:kaapa1}
\kappa=[\alpha_{max}+2(1+\alpha_{max})\frac{(L\gamma+\tau)}{\tau}]
\end{equation}
From lemma \ref{lemm1:gmopcgm}
\begin{equation}\label{mest5}
    \alpha_{min}||G_k||^2\leq -G_k^Tp_k\leq ||G_k||||p_k||.
\end{equation}
Therefore, we complete the proof by combining \eqref{mest4} and \eqref{mest5}.

\end{proof}

\begin{lemma}
\label{lemmA1}
    Suppose all the assumptions A1 and A2 hold, then \[\lim_{k\rightarrow \infty}\alpha_k||p_k||=0\]
\end{lemma}
    
\begin{proof}
    Beginning from the line search \eqref{eq:backtrack} $ G(z_k)^Tp_k\leq -\zeta \alpha_k||p_k||^2 \|G(z_k)\|$\\ Therefore 
    \[
    G(z_k)^T(x_k-z_k)=-\alpha_kG(x_k+\alpha_kp_k)^Tp_k
    \]
    \[
    G(z_k)^T(x_k-z_k)\geq \zeta \alpha_k^2||p_k||^2\|G(z_k)\|
    \]
     \begin{equation}\label{one}
         G(z_k)^T(x_k-z_k)\geq \zeta ||x_k-z_k||^2\|G(z_k)\|
     \end{equation}
 Now, we apply the monotonicity of $G$ and A1. Therefore there is $x^*\in \Gamma$ such that $G(x^*)=0$ \\
    \[
    G(z_k)^T(x_k-x^*)=G(z_k)^T(x_k-z_k+z_k-x^*)
    \]
    \[
    G(z_k)^T(x_k-x^*)=G(z_k)^T(x_k-z_k)+G(z_k)^T(z_k-x^*)
    \]
    \[
    G(z_k)^T(x_k-x^*)\geq G(z_k)^T(x_k-z_k)+G(x^*)^T(z_k-x^*)
    \]
     \begin{equation}\label{two}
          G(z_k)^T(x_k-x^*)\geq G(z_k)^T(x_k-z_k)
     \end{equation}
    
    This implies that
    \begin{equation}\label{three}
       G(z_k)^T(x_k-x^*)\geq \zeta ||x_k-z_k||^2 \|G(z_k)\|
    \end{equation}
    Applying the non-expansive property of the projection operator, we get
    \[
   || x_{k+1}-x^*||^2=||\Pi_{\Gamma}(x_k- \gamma\mu_kG(z_k))-x^*||^2,
    \]
    \[
   || x_{k+1}-x^*||^2\leq ||(x_k- \gamma\mu_kG(z_K))-x^*||^2,
    \]
     \[
   || x_{k+1}-x^*||^2\leq ||x_k-x^*||^2-2 \gamma\mu_kG(z_k)^T(x_k-x^*)+\gamma^2\mu_k^2||G(z_k)||^2.
    \]
   
    using \eqref{two}, we obtain
     \[
         || x_{k+1}-x^*||^2\leq ||x_k-x^*||^2- 2\gamma\mu_kG(z_k)^T(x_k-z_k)+\gamma^2\mu_k^2||G(z_k)||^2.
     \]
     But $\mu_k^2=\frac{(G(z_k)^T(x_k-z_k))^2}{||G(z_k)||^4}$
   \[
         || x_{k+1}-x^*||^2\leq ||x_k-x^*||^2- \gamma(2-\gamma)\frac{(G(z_k)^T(x_k-z_k))^2}{||G(z_k)||^2}.
     \]
     Using \eqref{one}, we obtain 
    
    \begin{equation}\label{five}
       || x_{k+1}-x^*||^2\leq ||x_k-x^*||^2- \gamma(2-\gamma)\zeta^2 ||x_k-z_k||^4.
    \end{equation}
    Implying that 
    \begin{equation}
       0\leq || x_{k+1}-x^*||^2\leq ||x_k-x^*||^2
    \end{equation}
 So, the sequence $\{\|x_k-x^*\|\}$  is bounded below and non-increasing, hence $\{x_k\}$ is convergent.

From \eqref{five} we have 
\[
   || x_{k}-x^*||^2\leq ||x_0-x^*||^2-\gamma(2-\gamma)\zeta^2\sum_{j=0}^k ||x_j-z_j||^4. 
\]
This means that 
\[
  \gamma(2-\gamma)\zeta^2\sum_{j=0}^k ||x_j-z_j||^4  \leq ||x_0-x^*||^2- || x_{k}-x^*||^2\leq ||x_0-x^*||^2,
\]
hence
\begin{equation}\label{six}
  \gamma(2-\gamma)\zeta^2\sum_{j=0}^{\infty} ||x_j-z_j||^4  \leq ||x_0-x^*||^2<\infty
\end{equation}

% For all $k$ and thus $\{x_k\}$ is bounded hence having a convergent sub-sequence. Using assumption A2, we have 
% \[
% ||G_k||=||G(x_k)-G(x^*)||\leq L||x_k-x^*||
% \]
% implying that
% \[
% ||G(x_k)||\leq L||x_0-x^*||
% \]
%  Therefore $G$ is bounded on $\Gamma$ due to its continuity.\\
%  Now using the line search \eqref{eq:backtrack}, we have 
%  \[
%  \sigma ||x_k-z_k||=\frac{\sigma ||x_k-z_k||^2}{||x_k-z_k||}
%  \]
%  By applying the monotonicity of $G$ and Cauchy-Schwartz inequality, we obtain
% \[
%  \sigma ||x_k-z_k||\leq \frac{G(x_k)^T(x_k-z_k)}{||x_k-z_k||}\leq L||x_0-x^*||
%  \]    
%  Since  $G(x_k)$ is bounded, therefore $\{z_k\}$ is also bounded.\\
%  Thus
%  \[
%  ||G(z_k)||-||G(x_k)||\leq ||G(z_k)-G(x_k)||\leq L||z_k-x_k||
%  \]

%  \[
%  \sigma ||x_k-z_k||\leq L||x_0-x^*||
%  \]
%  \[
%   ||x_k-z_k||\leq \frac{L||x_0-x^*||}{\sigma }
%  \]
%  Applying triangle inequality and Lipschitz continuity of $G$ here, we obtain
%  \[
%  ||G(z_k)||-||G(x_k)||\leq L^2\frac{||x_0-x^*||}{\sigma}
%  \]
%  \[
%  ||G(z_k)||\leq L^2\frac{||x_0-x^*||}{\sigma }+L^2||x_0-x^*||=M
%  \]
%  From \eqref{six}
%  \[
% || x_{k}-x^*||^2\leq ||x_0-x^*||^2-\frac{\sigma^2}{M^2}\sum_{j=0}^k ||x_j-z_j||^4 
%  \]
%  This implies that 
%  \[
%  \sum_{k=0}^{\infty} ||x_k-z_k||^4<\infty
%  \]
 This completes the proof.
\end{proof}
% \begin{remark}
%    \label{rem:trustregion}
%    Recall, in trust region methods, the search direction $p_k$  is constrained to lie within the trust region, typically represented by a ball:
% $\|p_k\| \leq \Delta_k$ where $\Delta_k$
%   is the radius of the trust region. This constraint ensures that the step size doesn't go too far in regions where the quadratic approximation may not be valid.
% \end{remark}



\begin{theorem}
\label{theorem:globalcovMOP}
    Suppose $x_k$ is generated by Algorithm \ref{alg:gmopcg}, then 
    \begin{equation}\label{eq:mgcc}
       \lim_{k\rightarrow \infty}\inf ||G_k||=0 
    \end{equation}  
\end{theorem}
\begin{proof}
    Suppose that there is some $\epsilon>0$ such that $||G_k||>\epsilon$ for all $k$. 
    Using \eqref{mest5} together with this, we obtain
    \begin{equation}
    \label{eq:condition}
    \alpha_{min}\epsilon\leq \|p_k\| ~~~~~~\forall k.
    \end{equation}
    From \eqref{eq:condition}, we have $\alpha_k\rightarrow 0$ as $k\rightarrow \infty$.\\
    Applying the line search in \eqref{eq:backtrack} there is $\bar{\alpha}_k$ (that is if $\alpha_k$ is the optimal step-size that satisfies the line search, then $\bar{\alpha}_k=\frac{\alpha_k}{\rho}$ violets the line search) such that 
    \begin{equation}
    \label{eq:violets}
    -G(x_k+\bar{\alpha}_k\zeta p_k)^Tp_k< \bar{\alpha}_k\zeta||p_k||^2\Pi_{[\zeta_1, \zeta_2]}(\|G(x_k + \bar{\alpha}_k p_k)\|)\leq \bar{\alpha}_k\zeta||p_k||^2\zeta_2
    \end{equation}
    Since $x_k$ and $p_k$ are bounded, we can select sub-sequences that converge to their accumulation points. Let $\bar{x}$ and $\bar{p}$ be the accumulation points of $x_k$ and $p_k$ respectively, then as $k$ approaches infinity and by continuity of $G$ in \eqref{eq:violets}, we have 
    \[
    -G(\bar{x})^T\bar{p}\leq 0
    \]
    Also using \eqref{eq:sufficientdescentmop} as $k$ approaches infinity, we have 
    \[
    -G(\bar{x})^T\bar{p}\geq 0
    \]

    This is for sure a contradiction.
    
    Therefore 
    \[
    \lim_{k\rightarrow \infty}\inf ||G(x_k)||=0.
    \]
    Hence global convergence. This completes the proof.
    
\end{proof}
We can provide an alternative proof that depends on the Lipschitz continuity of $G$. Assume $G$ is Lipschitz continuous.
\begin{proof}
We have two cases to consider.
\begin{enumerate}
    \item Case 1\\
    Suppose $\lim_{k\rightarrow \infty}\inf\|p_k\|=0$. Now when we apply \eqref{eq:sufficientdescentmop}, it means there is $\Omega \in \R_+$ such that  
    \[
    \|G_k\|\leq \Omega\|p_k\|, ~\forall ~k
    \]
    and taking the limits concludes \eqref{eq:mgcc}.
    \item Case II\\
    Let $\lim_{k\rightarrow \infty}\inf\|p_k\|\neq 0$.  Applying the line search in \eqref{eq:backtrack} there is $\Tilde{\alpha}_k$ (that is if $\alpha_k$ is the optimal step-size that satisfies the line search, then $\Tilde{\alpha}_k=\frac{\alpha_k}{\rho}$ violets the line search) such that 
    \begin{equation}
    \label{eq:violetsmopcgm}
    -G(x_k+\Tilde{\alpha}_k p_k)^Tp_k< \zeta\Tilde{\alpha}_k||p_k||^2\Pi_{[\zeta_1, \zeta_2]}(\|G(x_k + \Tilde{\alpha}_k p_k)\|)\leq \Tilde{\alpha}_k\zeta||p_k||^2\zeta_2
    \end{equation}
    But using sufficient descent condition \eqref{eq:mgcc}, Triangle inequality, Cauchy-Schwartz inequality and Lipschitz continuity of $G$, we have
    \[
    \alpha_{min} ||G_k||^2 \leq -G_k^Tp_k=(G(x_k+\Tilde{\alpha}_kp_k)-G_k)^Tp_k-G(x_k+\Tilde{\alpha}_kp_k)^Tp_k,
    \]
     \[
    \alpha_{min} ||G_k||^2 \leq L\Tilde{\alpha}_k||p_k||^2+\zeta\Tilde{\alpha}_k||p_k||^2\zeta_2,
    \]
     \[
    \alpha_{min} ||G_k||^2 \leq \Tilde{\alpha}_k||p_k||^2\left(L+\zeta\zeta_2\right).\]
    We arrive at
    \[
    \frac{\rho\alpha_{min}}{||p_k||(L+\zeta\zeta_2 )}\|G_k\|\leq \alpha_k||p_k||.
    \]
    Consequently, we have
    \[
     ||G_k||^2\leq \alpha_k||p_k||\frac{\alpha_{min}||p_k||(L+\zeta\zeta_2) }{\rho\alpha_{min}}.
    \]
    Taking the limits in $k$ implies that 
    \[
    \lim_{k\rightarrow \infty}\inf ||G_k||=0.
    \]
    Hence global convergence. This completes the proof.
\end{enumerate} 
\end{proof}

}
\subsection{Generalization of CGPM}
The CG parameter given by Hager and Zhang (HZ)~\cite{hager2005new} is a particular case of  Dai and Liao (DL) parameter ~\cite{dai2012another}. It happens when $t=2\frac{||y_{k-1}||^2}{s_{k-1}^Ty_{k-1}}$ in \eqref{eq:HZ} which may be considered to be an adaptive type of the D-L scheme.\\
We now propose a generalized parameter related to H-Z parameter $t_k$ to be 
\begin{equation}\label{eq:newhz}
t_k=\lambda \frac{||y_{k-1}||^2}{s_{k-1}^Ty_{k-1}},
\end{equation}
such that $\lambda > 0$. The proposed search direction is 
\begin{equation}
p_k=\begin{cases}\label{GCGPM}
    -G_k & k=0\\
    -\lambda_k G_k+\theta_k^{GCGM}p_{k-1}+ \tau a_kw_{k-1}& k\geq 1
\end{cases},
\end{equation}
such that $\lambda_k > 0$,
\begin{equation}
\label{eq:optimalcgpmpara}
    \theta_k^{GCGM}=\frac{G_k^Tw_{k-1}}{p_{k-1}^Tw_{k-1}}-\lambda_k\frac{||w_{k-1}||^2}{p_{k-1}^Tw_{k-1}}\frac{G_k^Tp_{k-1}}{p_{k-1}^Tw_{k-1}}, 
\end{equation}
where $\lambda_k=\Pi_{[\alpha_{min}, \alpha_{max}]}(\max(\frac{\|w_{k-1}\|^2}{s_{k-1}^Tw_{k-1}}, \frac{s_{k-1}^Tw_{k-1}}{\|s_{k-1}\|^2}))$ as used in \cite{sabi2020two} and 
\[
\Pi_{[\alpha_{min}, \alpha_{max}]}(x)=\max(\alpha_{min}, \min(x, \alpha_{max})),
\]

$a_k=\frac{G_k^Tp_{k-1}}{w_{k-1}^Tp_{k-1}}$, $\tau >0$, $w_{k-1}=y_{k-1}+r_kp_{k-1}$, $y_{k-1}=G(x_k)-G(x_{k-1})$, $s_{k-1}=x_k-x_{k-1}$ and  $r_k=1+\max\{0,-\frac{G_k^Tp_{k-1}}{w_{k-1}^Tp_{k-1}} \}$ \\
The next lemma verifies that \eqref{GCGPM} satisfies the sufficient descent condition. 


\begin{lemma}
\label{lemm1:gcgpm}
Let $\{p_k\}$ and $G(x_k)$ be produced by the algorithm \ref{algorithm:gcgpm}. Let  $0\leq \tau\leq 1$, and $\alpha_{min} \geq \frac{1+\tau}{2}$  , then  \eqref{GCGPM} meets the sufficient descent  condition $G_k^Tp_k\leq -\xi\|G_k\|^2$, where $\xi\geq 0$.
% $G_k^Tp_k\leq-\alpha_{min}(1-\frac{(1+\tau)^2}{4\alpha_{min}^2}) ||G_k||^2$

\end{lemma}

The condition holds for $k=0$.
\[
G_k^Tp_k=-\lambda_k ||G_k||^2+\theta_k^{GCGM}G_k^Tp_{k-1}+\tau a_kG_k^Tw_{k-1}.
\]
\[
G_k^Tp_k=-\lambda_k ||G_k||^2+\frac{G_k^Tw_{k-1}}{p_{k-1}^Tw_{k-1}}G_k^Tp_{k-1}-\lambda_k\frac{||w_{k-1}||^2}{p_{k-1}^Tw_{k-1}}\frac{G_k^Tp_{k-1}}{p_{k-1}^Tw_{k-1}}G_k^Tp_{k-1}+\tau \frac{G_k^Tp_{k-1}}{w_{k-1}^Tp_{k-1}}G_k^Tw_{k-1}.
\]
\[
G_k^Tp_k=-\lambda_k ||G_k||^2+(1+\tau)\frac{G_k^Tw_{k-1}}{p_{k-1}^Tw_{k-1}}G_k^Tp_{k-1}-\lambda_k\frac{||w_{k-1}||^2}{p_{k-1}^Tw_{k-1}}\frac{G_k^Tp_{k-1}}{p_{k-1}^Tw_{k-1}}G_k^Tp_{k-1}.
\]
\[
G_k^Tp_k=-\lambda_k ||G_k||^2+\frac{\sqrt{2\lambda_k}(1+\tau)}{\sqrt{2\lambda_k}}\frac{G_k^Tw_{k-1}}{p_{k-1}^Tw_{k-1}}G_k^Tp_{k-1}-\lambda_k\frac{||w_{k-1}||^2}{p_{k-1}^Tw_{k-1}}\frac{G_k^Tp_{k-1}}{p_{k-1}^Tw_{k-1}}G_k^Tp_{k-1}.
\]
\[
G_k^Tp_k=-\lambda_k ||G_k||^2+\frac{\sqrt{2\lambda_k}(1+\tau)}{\sqrt{2\lambda_k}}\frac{G_k^Tw_{k-1}G_k^Tp_{k-1}p_{k-1}^Tw_{k-1}}{(p_{k-1}^Tw_{k-1})^2}-\lambda_k\frac{||w_{k-1}||^2(G_k^Tp_{k-1})^2}{(p_{k-1}^Tw_{k-1})^2}.
\]
\[
G_k^Tp_k=-\lambda_k ||G_k||^2+\frac{\frac{(1+\tau)}{\sqrt{2\lambda_k}}G_k^T(p_{k-1}^Tw_{k-1})\sqrt{2\lambda_k}w_{k-1}(G_k^Tp_{k-1})}{(p_{k-1}^Tw_{k-1})^2}-\lambda_k\frac{||w_{k-1}||^2(G_k^Tp_{k-1})^2}{(p_{k-1}^Tw_{k-1})^2}.
\]
\[
G_k^Tp_k\leq -\lambda_k ||G_k||^2+\frac{\frac{(1+\tau)^2}{2\lambda_k}||G_k||^2(p_{k-1}^Tw_{k-1})^2+2\lambda_k
||w_{k-1}||^2(G_k^Tp_{k-1})^2}{2(p_{k-1}^Tw_{k-1})^2}-\lambda_k\frac{||w_{k-1}||^2(G_k^Tp_{k-1})^2}{(p_{k-1}^Tw_{k-1})^2}.
\]
\[
G_k^Tp_k\leq -\lambda_k ||G_k||^2+\frac{(1+\tau)^2}{4\lambda_k}||G_k||^2.
\]
But $0<\alpha_{min}\leq\lambda_k\leq \alpha_{max}$, then 
\begin{equation}
\label{eq:cgpmsuffientcondition}
    G_k^Tp_k\leq -\alpha_{min}(1-\frac{(1+\tau)^2}{4\alpha_{min}^2})||G_k||^2.
\end{equation}


This completes the proof.
\newpage
\subsubsection{Generalized CGPM Algorithm}
{\renewcommand{\baselinestretch}{0.95}
{\begin{algorithm}[H]
    \caption{Generalized CGPM}
    \label{algorithm:gcgpm}
    \begin{algorithmic}[1]
        \State \textbf{Input:} Function $G$, initial guess $x_0 \in \mathbb{R}^n$, 
        $\lambda_o > 0$, tolerance $\epsilon = \text{Tol}$, 
        $\rho \in (0, 1)$, parameters $\tau > 0$, $\eta > 0$, $\zeta > 0$, $\gamma, \gamma_1, \gamma_2, \gamma_3, \gamma_4 \in (0, 2) $, $0<\alpha_{min}\leq \alpha_{max}$, $0<\zeta_1\leq \zeta_2$,  projection function 
        $\Pi_{\Gamma}$ on convex set $\Gamma$
        \State \textbf{Output:} Solution $x^*$
        \State \textbf{Initialization:} Set $k \gets 0$, $z_k \gets x_0$, $\lambda \gets \lambda_o$, initialize $\alpha_k$
        
        \While{$\| G_k \| > \epsilon$}
        \If{$\|G(x_{k+1})\| < \|G(x_k)\|$ }
                \State $\lambda \gets \lambda$
                \State \textbf{Else}
                {$\lambda \gets \Pi_{[\alpha_{min}, \alpha_{max}]}(\max(\frac{\|w_{k-1}\|^2}{s_{k-1}^Tw_{k-1}}, \frac{s_{k-1}^Tw_{k-1}}{\|s_{k-1}\|^2}))$}
            \EndIf
            \State Compute $\theta_k^{GCGPM} $ using \eqref{eq:optimalcgpmpara}.
            \State Determine $p_k$ as in \eqref{GCGPM}
            \State Adjust $\alpha_k = \max \{ \rho^i \eta \mid i = 0, 1, 2, \dots \}$ such that 
            \begin{equation}
                G(x_k + \alpha_k p_k)^T p_k \leq -\zeta \alpha_k \| p_k \|^2 \Pi_{[\zeta_1, \zeta_2]}(\| G(x_k + \alpha_k p_k) \|)
                \label{eq:backtrack2}
            \end{equation}
            \If{$z_k = x_k + \alpha_k p_k \in \Gamma$ and $\| G(z_k) \| < \epsilon$}
                \State $x^* \gets z_k$
                \State \textbf{break}
            \EndIf
            \State Compute $\mu_k \gets \frac{G(z_k)^T (x_k - z_k)}{\| G(z_k) \|^2}$
            \State Update $x_{k+1} \gets \Pi_{\Gamma}(x_k - \gamma \mu_k G(z_k))$
            \State Compute $v_{k-1}=G(x_k)-G(x_{k-1})+r_kp_{k-1}$
             \If{$\|f_{k+1}\| < \|f_k\|$}
    \State $\gamma = \min(\gamma \cdot \gamma_1, \gamma_2)$
\Else
    \State $\gamma = \max(\gamma \cdot \gamma_3, \gamma_4)$
    \State \textbf{break}
\EndIf
            \If{$\|p_k\|  \approx 0$ }
                \State $x^* \gets x_k$
                \State \textbf{break}
            \EndIf
            \State Set $x_k \gets x_{k+1}$
        \EndWhile
        \State \textbf{Return} $x^*$
    \end{algorithmic}
\end{algorithm}
\renewcommand{\baselinestretch}{1.50}
}
\newpage
\subsubsection{Convergence analysis for the generalized CGPM}
\label{Convergence GCGPM}

{Assuming all the conditions A1 and A2 are met, then we can analyze the convergence of algorithm 1.
\begin{lemma}
    Let $G$ be Lipschitz continuous. Let  $0\leq \tau\leq 1$, and $\alpha_{min} \geq \frac{1+\tau}{2}$, Then 
    \[
    \alpha_k\geq \min\{\eta, \frac{\rho[4\alpha_{min}^2-(1+\tau)^2]}{4\alpha_{min}[L+\zeta\zeta_2]}\frac{||G_k||^2}{||p_k||^2}\}
    \]
\end{lemma}
\begin{proof}
   Let $\alpha_k$ be the optimal step length that satisfies \eqref{eq:backtrack2}, then, $\Tilde{\alpha}_k = \alpha_k/\rho $ violets \eqref{eq:backtrack2}. Therefore, 
\[
G(x_k+\tilde{\alpha}_kp_k)^Tp_k>- \zeta \tilde{\alpha}_k||p_k||^2\Pi_{[\zeta_1, \zeta_2]}(||G(x_k+\tilde{\alpha}_kp_k)||)
\]
But $G_k^Tp_k\leq -\alpha_{min}[1-\frac{(1+\tau)^2}{4\alpha_{min}^2}]||G_k||^2$. So
\[
||G_k||^2\leq -\frac{4\alpha_{min}}{[4\alpha_{min}^2-(1+\tau)^2]}G_k^Tp_k
\]
\[
||G_k||^2\leq \frac{4\alpha_{min}}{\left[4\alpha_{min}^2-(1+\tau)^2\right]}[(G(x_k+\Tilde{\alpha}_kp_k)- G_k)^Tp_k-G(x_k+\Tilde{\alpha}_kp_k)^Tp_k]
\]
\[
||G_k||^2\leq \frac{4\alpha_{min}}{\left[4\alpha_{min}^2-(1+\tau)^2\right]}[(G(x_k+\Tilde{\alpha}_kp_k)- G_k)^Tp_k+\zeta \tilde{\alpha}_k||p_k||^2\Pi_{[\zeta_1, \zeta_2]}(\|G(x_k+\tilde{\alpha}_kp_k)\|)],
\]

\[
||G_k||^2\leq \frac{4\alpha_k\alpha_{min}}{\rho[4\alpha_{min}^2-(1+\tau)^2]}[L+\zeta\zeta_2]||p_k||^2.
\]

\end{proof}
%\begin{lemma}
    %The search direction $p_k$ generated in algorithm 1 satisfies the trust region property.
%\end{lemma}
\begin{lemma}
\label{lemmacgpm}
    Suppose all the assumptions A1 and A2 hold, then \[\lim_{k\rightarrow \infty}\alpha_k||p_k||=0\]
\end{lemma}
    
\begin{proof}
    Beginning from the line search \eqref{eq:backtrack2} $ G(z_k)^Tp_k\leq -\zeta \alpha_k||p_k||^2||G(z_k)|| $\\ Therefore 
    \[
    G(z_k)^T(x_k-z_k)=-\alpha_kG(x_k+\alpha_kp_k)^Tp_k
    \]
    \[
    G(z_k)^T(x_k-z_k)\geq \zeta \alpha_k^2||p_k||^2||G(z_k)||
    \]
     \begin{equation}\label{cpart1}
         G(z_k)^T(x_k-z_k)\geq \zeta ||x_k-z_k||^2||G(z_k)||
     \end{equation}
 Now, we apply the monotonicity of $G$ and A1. Therefore there is $x^*\in \Gamma$ such that $G(x^*)=0$ \\
    \[
    G(z_k)^T(x_k-x^*)=G(z_k)^T(x_k-z_k+z_k-x^*)
    \]
    \[
    G(z_k)^T(x_k-x^*)=G(z_k)^T(x_k-z_k)+G(z_k)^T(z_k-x^*)
    \]
    \[
    G(z_k)^T(x_k-x^*)\geq G(z_k)^T(x_k-z_k)+G(x^*)^T(z_k-x^*)
    \]
     \begin{equation}\label{cpart2}
          G(z_k)^T(x_k-x^*)\geq G(z_k)^T(x_k-z_k)
     \end{equation}
    
    This implies that
    \begin{equation}
       G(z_k)^T(x_k-x^*)\geq \zeta ||x_k-z_k||^2||G(z_k)|| 
    \end{equation}
    Applying the non-expansive property of the projection operator, we get
    \[
   || x_{k+1}-x^*||^2=||\Pi_{\Gamma}(x_k- \gamma\mu_kG(z_k))-x^*||^2
    \]
    \[
   || x_{k+1}-x^*||^2\leq ||(x_k- \gamma\mu_kG(z_K))-x^*||^2
    \]
     \[
   || x_{k+1}-x^*||^2\leq ||x_k-x^*||^2-2 \gamma\mu_kG(z_k)^T(x_k-x^*)+\gamma^2\mu_k^2||G(z_k)||^2
    \]
     \begin{equation}\label{cconv1}
         || x_{k+1}-x^*||^2\leq ||x_k-x^*||^2- 2\gamma\mu_kG(z_k)^T(x_k-z_k)+\gamma^2\mu_k^2||G(z_k)||^2
     \end{equation}

     \begin{equation}
         || x_{k+1}-x^*||^2\leq ||x_k-x^*||^2- 2\gamma\frac{(G(z_k)^T(x_k-z_k))^2}{\|G(z_k)\|^2}+\gamma^2\frac{(G(z_k)^T(x_k-z_k))^2}{\|G(z_k)\|^2}
     \end{equation}
   
    \[
   || x_{k+1}-x^*||^2\leq ||x_k-x^*||^2-\gamma(2-\gamma) \frac{(G(z_k)^T(x_k-z_k))^2}{||G(z_k)||^2}
    \]
    \begin{equation}\label{cconv2}
        || x_{k+1}-x^*||^2\leq ||x_k-x^*||^2- \gamma(2-\gamma)(\zeta ||x_k-z_k||^4)
    \end{equation}
   
Inequalities \eqref{cconv1} and \eqref{cconv2} follow from \eqref{cpart1} and \eqref{cpart2}  respectively.\\
From \eqref{cconv2}, we obtain
\[
0\leq || x_{k+1}-x^*||^2\leq ||x_k-x^*||^2
\]
This implies that $\{||x_0-x^*||\}$ is a decreasing sequence and bounded below. This means that $\{x_k\}$ is convergent.
From \eqref{cconv2} we settle that 
\begin{equation}\label{cconv3}
   || x_{k}-x^*||^2\leq ||x_0-x^*||^2- \gamma(2-\gamma)\zeta^2\sum_{j=0}^k( ||x_j-z_j||^4.
\end{equation}

 
 \[
 \sum_{k=0}^{\infty} ||x_k-z_k||^4\leq ||x_0-x^*||^2 <\infty.
 \]
 This completes the proof.
\end{proof}
%  \begin{lemma}
%  Let $\R^n$ be the solution set of \eqref{eq:op} and $s_{k-1}=x_k-x_{k-1}$. If the search direction $p_k$ be generated by  Algorithm \ref{algorithm:gcgpm}, then $p_k$ satisfies the trust region property
%      \[
%     \alpha_{min}(1-\frac{(1+\tau)^2}{4\alpha_{min}^2}) \|G(x_k)\|\leq \|p_k\|\leq \kappa \|G(x_k)\|.
%      \]
    
%  \end{lemma}
%  \begin{proof}
%  From $\Gamma=\R^n$, then $\Pi_{\Gamma}(x_{k-1}-\gamma\mu_{k-1}G(z_{k-1}))=x_{k-1}-\gamma\mu_{k-1}G(z_{k-1})=x_k$
%  But 
%  \[
%  p_k=-\lambda_kG_k+\theta_k^{GCGPM}p_{k-1}+\tau a_kw_{k-1},
%  \]
% where $\theta_k^{GCGPM}=\frac{w_{k-1}^TG_k}{w_{k-1}^Tp_{k-1}}-\lambda_k\frac{\|w_{k-1}\|^2}{w_{k-1}^Tp_{k-1}}\frac{p_{k-1}^TG_k}{w_{k-1}^Tp_{k-1}}$, $a_k=\frac{G_k^Tp_{k-1}}{w_{k-1}^Tp_{k-1}}$ and $w_{k-1}=G_k-G_{k-1}+r_kp_{k-1}$.
% So
% \[
% \|w_{k-1}\|\leq L\|s_{k-1}\|+r_k\|s_{k-1}\|.
% \]
% Also we have that $r_k=1+\max\{0, -\frac{y_{k-1}^Tp_{k-1}}{\|G_k\|^2}\}=1$ because $G$ is monotone.\\
% This means that 
% \[
% \|w_{k-1}\|\leq (L+1)\|s_{k-1}\|.
% \]
% Again
% \[
% w_{k-1}^Tp_{k-1}=y_{k-1}^Tp_{k-1}+r_ks_{k-1}^Tp_{k-1}
% \]
% This implies that 
% \[
% w_{k-1}^Tp_{k-1}\geq s_{k-1}^Tp_{k-1}=\gamma\alpha_{k-1}\|p_
% \]
%  \end{proof}

\begin{theorem}
\label{theorem:globalconvcgpm}
    Suppose $x_k$ is produced by  Algorithm \ref{algorithm:gcgpm} and $\alpha_{min}> \frac{(1+\tau)}{2}$ as defined in algorithm \ref{algorithm:gcgpm}, then 
    \begin{equation}\label{cgcc}
       \lim_{k\rightarrow \infty}\inf ||G_k||=0 .
    \end{equation}  
\end{theorem}
\begin{proof}
We have two cases to consider.
\begin{enumerate}
    \item Case 1\\
    Suppose $\lim_{k\rightarrow \infty}\inf\|p_k\|=0$. Now when we apply \eqref{eq:cgpmsuffientcondition}, it means there is $\Omega \in \R_+$ such that  
    \[
    \|G_k\|\leq \Omega\|p_k\|, ~\forall ~k
    \]
    and taking the limits concludes \eqref{cgcc}.
    \item Case II\\
    Let $\lim_{k\rightarrow \infty}\inf\|p_k\|\neq 0$.  Applying the line search in \eqref{eq:backtrack2} there is $\Tilde{\alpha}_k$ (that is if $\alpha_k$ is the optimal step-size that satisfies the line search, then $\Tilde{\alpha}_k=\frac{\alpha_k}{\rho}$ violets the line search) such that 
    \begin{equation}
    \label{eq:violetscgpm}
    -G(x_k+\Tilde{\alpha}_k p_k)^Tp_k< \zeta\Tilde{\alpha}_k||p_k||^2\Pi_{[\zeta_1, \zeta_2]}(\|G(x_k + \Tilde{\alpha}_k p_k)\|)\leq \Tilde{\alpha}_k\zeta||p_k||^2\zeta_2
    \end{equation}
    But using sufficient descent condition \eqref{cgcc}, Triangle inequality, Cauchy-Schwartz inequality and Lipschitz continuity of $G$, we have
    \[
    \alpha_{min}[1-\frac{(1+\tau)^2}{4\alpha_{min}^2}] ||G_k||^2 \leq -G_k^Tp_k=(G(x_k+\Tilde{\alpha}_kp_k)-G_k)^Tp_k-G(x_k+\Tilde{\alpha}_kp_k)^Tp_k,
    \]
     \[
    \alpha_{min}[1-\frac{(1+\tau)^2}{4\alpha_{min}^2}] ||G_k||^2 \leq L\Tilde{\alpha}_k||p_k||^2+\zeta\Tilde{\alpha}_k||p_k||^2\zeta_2,
    \]
     \[
    \alpha_{min}[1-\frac{(1+\tau)^2}{4\alpha_{min}^2}] ||G_k||^2 \leq \Tilde{\alpha}_k||p_k||^2\left(L+\zeta\zeta_2\right).\]
    We arrive at
    \[
    \frac{\rho [4\alpha_{min}^2-(1+\tau)^2] }{4||p_k||\alpha_{min}(L+\zeta\zeta_2 )}\|G_k\|\leq \alpha_k||p_k||.
    \]
    Consequently, we have
    \[
     ||G_k||^2\leq \alpha_k||p_k||\frac{4\alpha_{min}||p_k||(L+\zeta\zeta_2) }{\rho[4\alpha_{min}^2-(1+\tau)^2]}.
    \]
    Taking the limits in $k$ implies that 
    \[
    \lim_{k\rightarrow \infty}\inf ||G_k||=0.
    \]
    Hence global convergence. This completes the proof.
\end{enumerate} 
\end{proof}
}
\section{Numerical Experiments}
\label{testing}
{In this section, we studied and compared the performance of the GCGPM and GMOPCGM with the other three methods. That is STTDFPM [\cite{ibrahim2023two}, Algorithm 1], MOPCGM [\cite{sabi2023modified}, Algorithm 2.1], and CGPM [\cite{zheng2020conjugate}, Algorithm 2.1] without changing their line searches or parameters.\\  The comparison was conducted by considering the number of iterations, CPU time, and function evaluation number.\\ The 
 Algorithms were applied to Julia and tested with the same initial values. 


 
 % \begin{align}
 %      \begin{multlined}[t]
 %        x_0^1=(0,\cdots, 0)^T,  x_0^2=(0.2, \cdots, 0.2)^T, x_0^3=(0.4, \cdots, 0.4)^T ,  x_0^4=(0.6, \cdots, 0.6)^T \\x_0^5=(0.8, \cdots, 0.8)^T,  x_0^6=(1, \cdots, 1)^T, x_0^7=(1.1, \cdots, 1.1)^T
 %        x_0^8=(8, 8 \cdots, 8)^T\\
 %        %x_0^9=(1,\frac{1}{2}, \frac{1}{3}, \cdots, \frac{1}{n})^T,  x_0^{10}=(0,\frac{1}{n}, \frac{2}{n} \cdots, \frac{n-1}{n} )^T, x_0^{11}=(\frac{1}{n}, \cdots, \frac{1}{n})^T, x_0^{12}=(\frac{1}{2}, \frac{1}{2^2} \cdots, \frac{1}{2^n})^T
 %      \end{multlined}
 % \end{align}

 The number of iterations, function evaluations, and CPU time were represented by IT, FE, and CPU, respectively, and the norm of $G$ was also recorded to notice the accuracy of the algorithms. 19 test problems from different sources were used. All the experiments were conducted with three distinct dimensions, that is $10^3$, $10^4$, and $5\times 10^4$. 14 different initial points were used for each dimension in the experiment.\\
 We carefully selected the best parameters that seemed to suit each proposed algorithm for better performance.

    \begin{figure}[h] % Place image here
    \centering
    \includegraphics[width=0.8\textwidth]{Scatter_F14-GMOP.png} % Adjust width as needed
    \caption{Scatter diagrams for parameter values for GMOPCGM}
    \label{fig:scater diagram1} % Reference label
\end{figure}


    \begin{figure}[h] % Place image here
    \centering
    \includegraphics[width=0.8\textwidth]{Scatter_F15-GCGPM.png} % Adjust width as needed
    \caption{Scatter diagrams for parameter values for GCGPM}
    \label{fig:scater diagram2} % Reference label
\end{figure}
% Figures \ref{fig:scater diagram1} and \ref{fig:scater diagram2} show how the patterns of $\tau=\sigma$  and $\lambda_o=\phi$ for both GMOPCGM and GCGPM. We studied these two parameters and realized the following.
% \begin{enumerate}
%     \item GMOPCGM
%     \begin{enumerate}
%         \item[(i)]$\tau$ did not have much impact on the algorithm's performance. Due to this, we selected $\tau=1.0$
%         \item[(ii)] The performance of the algorithm improved when $\lambda_o=\phi$ was smaller than $2$. Because of this, we decided $\lambda_o$ to  take the values in  $(0, 2]$   
%     \end{enumerate}
%     \item GCGPM
%     \begin{enumerate}
%         \item[(i)] $\tau$ significantly affected the  performance of the algorithm when it was near zero than any other for almost all the test problems. This led us to choose $\tau=0.001$.
%         \item[(ii)] The performance of the algorithm further improved when $\lambda_o=\phi$ was smaller than $2$ but more than $0.5$.due to  this, we decided $\lambda_o$ to  take the values in $(0.5, 2]$
%     \end{enumerate}
% \end{enumerate}
In order to study the effect of parameters $\tau$ and $\phi$, we utilized machine learning-based hyperparameter selection approaches. Specifically, we have utilized the Bayesian Optimization (BO) technique to explore and exploit the parameter space. The key idea in BO is to learn a surrogate model that captures the optimization approach. We have considered a Tree-structured Parzen Estimator (TPE) approach, where the kernel density is estimated to learn dynamics between the parameter space and objective function. In addition to that, we have used non-surrogate approaches like Random Sampling (RS) and Grid Search (GS) to identify the values of the two parameters to explore and evaluate.
 
In all three approaches, both parameters are searched in the following intervals: $[0.1,10]$ for GMOPCGM and GCGPM. Moreover, we have performed 50 iterations of the two-parameter evaluations, where each iteration involves a combination of the two parameters. The study was carried out on different objective functions. Within each objective function, over multiple dimensions, we analyzed the relationship between the objective function value and the parameter values.   
 
For example, in Figure \ref{fig:scater diagram1}, results of the study on F14 using GMOPCGM are depicted. The three columns of scatter plots correspond to TPE, random, and grid search-based sampling, respectively. In a given column, the five scatter plots corresponding to $n=10,50, 100, 500, 1000$, respectively, are illustrated. Similarly, in Figure \ref{fig:scater diagram2}, results of the study on F15 using GCGPM are presented. Based on similar studies conducted over different functions w.r.t GMOPCGM \& GCGPM, we have observed the following:
\begin{enumerate}
	\item GMOPCGM
	\begin{enumerate}
		\item[(i)]$\tau$ did not have much impact on the algorithm's performance. Due to this, we selected $\tau=1.0$
		\item[(ii)] In almost all the cases, the performance of the algorithm improved when $\lambda_o=\phi$ was around $1.5$. Because of this, we decided $\lambda_o$ to take the values in $(0, 2)$.
	\end{enumerate}
	\item GCGPM
	\begin{enumerate}
		\item[(i)] $\tau$ significantly improved the algorithm's performance when it was near zero for almost all the test problems. This led us to choose $\tau=0.001$.
		\item[(ii)] The performance of the algorithm was further improved when $\lambda_o=\phi$ was greater than $1$. Due to this, we decided $\lambda_o$ to take the values in $(0.5, 2]$.
	\end{enumerate}
\end{enumerate}
 For GCGPM, we selected the following parameters
$\tau=0.001$, $\eta=0.6$, $\lambda_o=1.0$, $\rho=0.5$, $\zeta=0.1$, $\zeta_1=1.0$, $\zeta_2=1.0$, $\alpha_{min}=0.55$, $\alpha_{max}=4.9$, $\gamma=1.8$, 
$\gamma_1=1.1$, $\gamma_2=1.7$,   $\gamma_3=1.05$, $\gamma_4=1.05$\\
 For GMOPCGM, we selected the following parameters\\
 $\tau=1.0$, $\rho=0.8$,  $\beta=0.5$,  $\zeta=0.0001$, $\alpha_{min}=0.1$, $\alpha_{max}=2.0$, $\lambda_o=1.0$, 
 $\gamma=1.1$, $\gamma_2=1.8$, $\gamma_3=1.0$, $\gamma_4=1.0$, $\zeta_1=1.0$, $\zeta_2=1.0$\\
 The following are the initial points used for the experiments;\\
 $0=(0,\cdots, 0)^T$, $0.2=(8.2, \cdots, 8.2)^T$, $0.4=(\frac{2}{5},\cdots, \frac{2}{5})^T$, $0.5=(\frac{1}{2}, \cdots, \frac{1}{2})^T$, $0.6=(\frac{3}{5}, \cdots, \frac{3}{5})^T$, $0.8=(\frac{4}{5}, \cdots, \frac{4}{5})^T$, $1.0=(1.0, \cdots, 1.0)^T$, $1.1=(\frac{11}{10}, \cdots, \frac{11}{10})^T$, $1-1/m=(1-\frac{1}{m}, \cdots, 1-\frac{1}{m})^T$, $1/m=(1, \frac{1}{2}, \frac{1}{3}, \cdots, \frac{1}{m})^T$, $(k-1)/m=(0, \frac{1}{m}, \frac{2}{m}, \cdots, \frac{m-1}{m})^T$, $1/m=(\frac{1}{m}, \frac{1}{m}, \cdots, \frac{1}{m})^T$, $1/3^k=(\frac{1}{3}, \frac{1}{3^2}, \cdots, \frac{1}{3^m})^T$, $k/m=(\frac{1}{m}, \frac{2}{m}, \cdots, 1)^T$.\\
 All algorithms were terminated if either of the following was met;
 \begin{enumerate}
     \item $\|G(x_k)\|<Tol$ \item $ p_k<0.1 Tol$ \item  $k>2000$
 \end{enumerate}
  $\|G(x_k)\|$  represents the usual norm of $G$ at $x_k$ and $Tol=10^{-11}$. \\
  For a better assessment and comparison of the performance of the various schemes, we employed the  Moré and Dolan performance profile in ~\cite{dolan2002benchmarking}. The performance profile of the three new algorithms and their counterparts are represented in Figures \ref{fig:functionevals}-\ref{fig:time}. As in the figures, the vertical axes depict the chances that a certain solver outperforms the rest of the competing algorithms.
\newpage
    \begin{figure}[h] % Place image here
    \centering
    \includegraphics[width=0.8\textwidth]{Function_Evaluations_2025_03_12_08_56.png} % Adjust width as needed
    \caption{Profile of function evaluations}
    \label{fig:functionevals} % Reference label
\end{figure}

    \begin{figure}[h] % Place image here
    \centering
    \includegraphics[width=0.8\textwidth]{Iterations_2025_03_12_08_56.png} % Adjust width as needed
    \caption{Profile of iterations}
    \label{fig:iterations} % Reference label
\end{figure}
\newpage
    \begin{figure}[h] % Place image here
    \centering
    \includegraphics[width=0.8\textwidth]{Time_2025_03_12_08_56.png} % Adjust width as needed
    \caption{Profile of time}
    \label{fig:time} % Reference label
\end{figure}
From Figures \ref{fig:functionevals}-\ref{fig:time}, it was observed that GCGPM outperformed its counterparts where it won by  $42\%$ $51\%$, and $41\%$ as depicted in Figures \ref{fig:functionevals}, \ref{fig:iterations} and \ref{fig:time} respectively. This is because the four methods, GCGPM, STTDFPM, and GMOPCGM, share the same properties, so this could likely cause a small deviation in their performance. This can be witnessed from the figures. However, There was a great improvement of the two generalized methods GCGPM and GMOPCGM compared to CGPM and MOPCGM respectively. In addition the GMOPCGM and STTDFPM were to close to each other in terms of Iterations and CPU time though the later slightly lagged behind in function evaluations. \\
The following are the problems used in the experiment.\\
{\bf Problem 5.1}(Problem 4.1 in ~\cite{sabi2020two}) $G(x)$ is \\
 $G(x_i) = 2x_i - \sin x_i, ~~i = 1, 2,\cdots, n$ and $\Gamma = [-2, \infty].$\\
{\bf Problem 5.2}(Problem 10 in ~\cite{la2006spectral}) \\
 $G(x_i) = \log(|x_i| + 1) - \frac{x_i}{n}, ~~
 , i = 1, 2, \cdots, n.$ and $\Gamma = \R.$\\
 {\bf Problem 5.3}(Problem 4.1 in\cite{zheng2020conjugate})\\
 $G(x_i) = \exp(x_i) - 1, ~~
 , i = 1, 2, \cdots, n.$ and $\Gamma = \R.$\\
 %{\bf Problem 5.3}\cite{la2006spectral} The general interpretation $G(x)$ defined as\\
 %$G(x_i) = \frac{i}{10}(1-x_i^2-\exp(x_i^2)), ~~
 %, i = 1, 2, \cdots, n-1.\\
 %G(x_n)=\frac{n}{10}(1-\exp(x_n^2))$\\
 {\bf Problem 5.4}(Problem 4.5 ~\cite{sabi2020two} )The general interpretation $G(x)$ defined as\\
 $G(x_i) = 4x_i+(x_{i+1}-2x_i)-\frac{x_{i-1}^2}{3}, ~~
 , i = 1, 2, \cdots, n-1.\\
 G(x_n)=4x_n+(x_{n-1}-2x_n)-\frac{x_{n-1}^2}{3}$ \\and $\Gamma = \R.$\\
 {\bf Problem 4.5}(Problem 4.4 in \cite{zheng2020conjugate}) Exponential problem $G(x)$ defined as\\
 $G(x_1) = x_1-\exp{\cos(\frac{x_1+x_2}{n+1})},\\
 G(x_i) = x_i-\exp{\cos(\frac{x_{i-1}+x_i+x_{i+1}}{n+1})}, ~~ i = 2, \cdots, n-1$.\\
 $ G(x_n) = x_n-\exp{\cos(\frac{x_{n-1}+x_n}{n+1})}$\\
 and $\Gamma =\R$\\
{\bf Problem 5.6}(Problem 4.4~\cite{zheng2020conjugate}) Exponential problem $G(x)$ defined as\\
 $G(x_1) = x_1+\sin(x_1)-1,\\
 G(x_i) = -x_{i-1}+2x_i+\sin(x_i)-1,~~ i = 2, \cdots, n-1.\\
 G(x_n) = x_n+\sin(x_n)-1$\\
 {\bf Problem 5.7}(Problem 19 in ~\cite{la2006spectral})Zero Jacobian function $G(x)$ defined as\\$
 G(x_1)=\sum_{j=1}^nx_j^2\\
 G(x_{i}) =  -2x_1x_i,~~~for ~~i=2,\cdots, n$\\ and 
 $\Gamma =\R$\\
{\bf Problem 5.8}(Problem 14~\cite{song2024efficient}) The general interpretation of $G(x)$ defined as\\$
  G(x_1) =x_1 (x_1^2 + x_2^2) - 1\\
  G(x_i) =  x_i  (x_{i-1}^2 + 2x_i^2 + x_{i+1}^2) - 1,~~~for ~~i=2,\cdots, n-1\\
  G(x_n) = x_n  (x_{n-1}^2 + x_n^2)$\\ and 
 $\Gamma =\R$\\
 {\bf Problem 5.9}(Problem 12~\cite{la2006spectral})Trigexp function  $G(x)$ defined as\\$
  G(x_1) =3x_1^3 + 2x_2 - 5 + \sin(|x_1-x_2|)\sin(|x_1+x_2|)\\
  G(x_i) =-x_{i-1}\exp(x_{i-1}-x_i) + x_1(4+3x_i^3) + 2x_{i+1} - 5 + \sin(|x_i-x_{i+1}|)\sin(|x_i+x_{i+1}|),~~~for ~~i=2,\cdots, n-1\\
  G(x_n) =  -x_{n-1}  \exp(x_{n-1} - x_n) + 4x_n - 3$\\ and 
 $\Gamma =\R$\\
 {\bf Problem 5.10}(Problem 2~\cite{song2024efficient}) Complementary problem
 $G(x)$ defined as\\$
 G(x_i) = (x_i-1)^2 - 1.01, ~~for ~i=1, \cdots, n$\\and 
 $\Gamma =\R$\\
 {\bf Problem 5.11}(Problem 4~\cite{song2024efficient}) Complementary problem and 
 $G(x)$ defined as\\$
 G(x_i) =\frac{i}{n} \exp{x_i}-1 , ~~for ~i=1, \cdots, n$\\and 
 $\Gamma =\R$\\
{\bf Problem 5.12}(Problem 11~\cite{ibrahim2024two})\\
 $ G(x_i) =x_i- \sin(|x_i - 1|) , ~~for ~i=1, \cdots, n$\\and 
 $\Gamma =\R$\\
{\bf Problem 5.13}(Problem 4.5 in ~\cite{waziri2022two})\\$
 G(x_i) =2x_i- \sin(|x_i - 1|) , ~~for ~i=1, \cdots, n$\\and 
 $\Gamma =\R$\\
{\bf Problem 5.14}(Problem 6~\cite{song2024efficient}) 
 \\$
 G(x_i) =x_i- 2\sin(|x_i - 1|) , ~~for ~i=1, \cdots, n$\\and 
 $\Gamma =\R$\\
{\bf Problem 5.15}(Problem 11~\cite{song2024efficient}) 
 \\$
 G(x_i) =(\exp{x_i})^2+3\sin x_i\cos x_i-1 , ~~for ~i=1, \cdots, n$\\and 
 $\Gamma =\R$\\
 {\bf Problem 5.16}(Problem 5~\cite{waziri2020descent}) The  singular function
 $G(x)$ defined as\\
 $G(x_1)=2.5x_1 +x_2 - 1\\
 G(x_i) = x_{i-1} + 2.5x_i + x_{i+1} - 1, ~~for ~i=2, \cdots, n-1\\
 G(x_n)=x_{n-1} + 2.5x_n - 1$\\
 {\bf Problem 5.17}(Problem 1~\cite{zhou2007limited}) $G(x)$ defined as\\$
 G(x_i) =  2x_i - \sin(|x_i|), ~~for ~i=1, \cdots, n$\\
 {\bf Problem 5.18}(Problem 32~\cite{la2006spectral}) Minimal function  $G(x)$ is defined as\\$
 G(x_i) = 0.5\{\log{x_i}+\exp{x_i} -\sqrt{(\log{x_i}-\exp{x_i})^2-10^{-10} }\}, ~~for ~i=1, \cdots, n$\\
 {\bf Problem 5.19}(Problem 4.11~\cite{li2021scaled}) $G(x)$ defined by\\$
 G(x_i) = 2(10^{-5})(x_i-1) + 4x_i \sum_{j=1}^nx_j^2 - x_i, ~~for ~i=1, \cdots, n$\\
{\bf Problem 5.20}(Problem 4.6~\cite{sabi2023modified})  $G(x)$ defined by\\$
 G(x_i) = x_{i}(\cos(x_i-1/n)(\sin x_i-1-(1-x_i)^2-1/n\sum_{j=1}^nx_j), ~~for ~i=1, \cdots, n$
 
 
\subsection{Signal Restoration}
{Signal restoration refers to recovering an original signal from degraded observed signals\cite{chen2024manifold}. Signal restoration is a real-world problem that includes but is not limited to dequantization ~\cite{liu2018graph, liu2016random},  denoising ~\cite{dinesh2020point, pang2017graph, zeng20193d}, deblurring~\cite{bai2018graph, chen2021fast}. \\
Signal restoration includes large-scale inverse problems in which a multidimensional signal $x$ is to be obtained from the observation of
data $y$ consisting of signals.  Both the original signal $x$ and the observed $y$ are taken to lie in some real Hilbert spaces which may be independent~\cite{combettes2005signal, stark2013image}.\\
The observed signal is given by 
\begin{equation}\label{eq:sr}
    y=Hx+k
\end{equation}such that $x$ is the signal that we want to recover from $y$, k is called the additive noise, $H$ is a given operator representing the observation process like blurring or degradation, ~\cite{selesnick2009sparse,soussen2011bernoulli}. 
The problem of restoring the original signal $x$ from the observed signal $y$ is an inverse problem ~\cite{espanol2010multilevel, selesnick2009sparse}. 

Different techniques have been developed to resolve \eqref{eq:sr}. For example, Multilevel approach~\cite{espanol2010multilevel}, majorization-minimization\cite{selesnick2009sparse}, Accelerated Projected Gradient Method  APGM\cite{daubechies2008accelerated}  and many more. Therefore \eqref{eq:sr} can be written as
\begin{equation}\label{eq:sp}
    T(x)= ||y-Hx||_2^2+\eta||x||_1
\end{equation}
so to obtain $x$, its through minimizing/possibly solving \eqref{eq:sp} called $l_1$ regularized linear inverse problem or the  penalized least-squares problem.\\
The presence of the $l_1$ term may lead to small components
of $x$ to become exactly zero, thus leading to sparse solutions
\cite{figueiredo2007gradient}.\\
Yin et al in \cite{yin2023family}, evaluated the performance of their technique by obtaining the sparse solution $x$ for 
\begin{equation}\label{eq:23}
    \min_{x\in \R^n} \frac{1}{2}||y-Hx||^2+\eta||x||_1
\end{equation}
and \eqref{eq:23} is assumed to contain under-determined linear systems of equations such that $H\in \R^{m\times n}$ and $y\in \R^m$ and $m<<n$.

Problem \eqref{eq:23} can be redefined by  
\begin{equation}\label{eq:24}
    \min \frac{1}{2}z^T Qz+ d^Tz,~~ z\geq 0
\end{equation}
such that $z=\begin{pmatrix}
\max\{x, 0\}\\ \max\{-x, 0\}
\end{pmatrix}$, $d=\begin{pmatrix}
\eta E_n-H^Ty\\ \eta E_n+H^T y
\end{pmatrix}$, $Q=\begin{pmatrix}
H^TH&-H^TH\\ -H^TH&H^TH
\end{pmatrix}$.\\ From \cite{gao2018efficient, xiao2011non} the function $G$ is monotone and continuous.\\It was also noted in ~\cite{sabi2020two} that $z$ satisfies \eqref{eq:24} if and only if 
it satisfies \eqref{eq:application}
\begin{equation}
    \label{eq:application}
  G(z)=\min\{z,  Qz+d\}=0
\end{equation}
The 'min' is interpreted to be a point-wise minimum and $G$ is monotone and continuous as proved in ~\cite{iusem1997newton}.\\Therefore problem \eqref{eq:application} can be interpreted to be in the form of \eqref{eq:op}. In the experiment, firstly, we defined a sparse signal $x_{original}$ of length $n=2^{12}$ and sparsity $k=2^{9}$, that is $k$ are non-zero elements in the signal that are randomly selected.  \\
A Sensing Matrix $H$ of size $m\times n$ was obtained randomly. Where $m<n$ and in this case we considered $m=2^{11}$.  The Gaussian noise was also determined whose components were produced normal distributions $N(\mu=0, \sigma=0.01)$. The observation $y$ was considered to be the sum of the noise and the product of the sensing matrix and the original signal. 
\newpage

    \begin{figure}[h] % Place image here
    \centering
    \includegraphics[width=0.8\textwidth]{reconstructed_signalsfgff.png} % Adjust width as needed
    \caption{Reconstructed signals}
    \label{fig:reconstructedsignals} % Reference label
\end{figure}

\newpage
In the application of the GCGPM and GMOPCGM in compressed sensing, the iteration process began with the initial point $x_o=H^Ty$ of length $n$, and the experiment was terminated when $\|G(x_k)\|< 10^{-5}$.\\
The parameters for the GCGPM remained the same.\\
For GMOPCGM, we selected the following parameters.\\
 $\tau=1.05$, $\rho=0.8$,  $\beta=0.5$,  $\zeta=0.0001$, $\alpha_{min}=0.1$, $\alpha_{max}=2.0$, $\lambda_o=1.0$, 
 $\gamma=1.1$, $\gamma_2=1.8$, $\gamma_3=0.85$, $\gamma_4=1.0$, $\zeta_1=1.0$, $\zeta_2=1.0$.


To consolidate the applicability of our schemes in recovering sparse signals, we compared them with STTDFPM in ~\cite{ibrahim2023two}, CGPM in ~\cite{zheng2020conjugate} and MOPCGM in ~\cite{sabi2023modified}. The comparison helped us to evaluate how our methods perform best relative to their counterparts. A good performance implies a small mean square error given by,
\[
MSE=\frac{\|x_{original}-x_{recovered}\|}{n}.
\]
So, for fairness to all solvers, we conducted 10 experiments and determined the means for the  Iterations, means for the function, and mean times besides the mean square errors. The results were recorded in table \ref{tab:algorithm_performance} below.
{\small 
\begin{table}[H]
    \centering
    \sisetup{scientific-notation=true, round-mode=places, round-precision=2}
    \begin{tabular}{lcccc}
        \toprule
        \textbf{Algorithm} & \textbf{Iterations} & \textbf{Function Evals } & \textbf{Time} & \textbf{MSE} \\
        \midrule
        MOPCM        & 421.7           & 1267.1 & 7.5477 & \num{5.22894E-8} \\
        GMOPCGM      & \textbf{269.0}  & \textbf{808.8}  & \textbf{4.5182} & \num{5.22928E-8} \\
        CGPM         & 1708.4          & 5468.7 & 39.6073 & \num{5.05665E-8} \\
        GCGPM        & 363.2           & 1089.6 & 8.0243 & \num{5.22875E-8} \\
        STTDFPM      & 398.2           & 831.7  & 6.8776 & \num{5.14822E-8} \\
        \bottomrule
    \end{tabular}
    \caption{Comparison of algorithm performance based on iterations, function evaluations, time, and MSE.}
    \label{tab:algorithm_performance}
\end{table}
}
As we can see from the table, it is clear that CGPM gives a small mean square error. However, GMOPCGM took less time, had few iterations, and had few function evaluations. Figure \ref{fig:reconstructedsignals} shows the quality of the reconstructed signals by the solvers.


}
\section*{Conclusion}
In this paper, we proposed two efficient methods, and GCGPM outperformed the rest of its counterparts. They satisfy the sufficient descent condition. They also converge globally under the Lipschitz continuity of the function.

\subsection{Data availability}
The data used to support the findings of this study are included in the article.
\subsection{Conflicts of Interest}
The authors declare that they have no conflicts of interest.
\newpage
%\printbibliography
%%%%%%%%%%%%%%%%%%%%%%%%%%%%%%%%%%%%%%%%%%%%%%%%%%%%%%%%%%%%%%%%%%%%%%%%%%%%%%%%%%%%%%%%%%%%%%%%%%%%%%%%%%%
%%%%%%%%%%%%%%%%%%%%%%%%%%%%%%%%%%%%%%%%%%%%%%%%%%%%%%%%%%%%%%%%%%%%%%%%%%%%%%%%%%%%%%%%%%%%%%%%%%%%%%%%%%%

%\newpage

%%%%%%%%%%%%%%%%%%%%%%%%%%%%%%%%%%%%%%%%%%%%%%%%%%%%%%%%%%%%%%%%%%%%%%%%%%%%%%%%%%%%%%%%%%%%%%%%%%%%%%%%%%%

\newpage

% \chapter*{Bibliography} % Unnumbered Chapter
\phantomsection % This ensures correct numbering in the table of contents
\addcontentsline{toc}{chapter}{References}%Add Acknowledgements to the table of contents as a clickable link with hyperref package
\label{references}
% \addtocontents{toc}{\contentsline{chapter}{\numberline{}REFERENCES}{\pageref{references}}}\label{references}
% \bibliographystyle{plain}
% \bibliography{library}
% \phantomsection % For hyperref
% \addcontentsline{toc}{chapter}{Bibliography} % Add to TOC
%\bibliographystyle{plain} % Choose your style
%\bibliography{library} % Point to your .bib file
\printbibliography


%%%%%%%%%%%%%%%%%%%%%%%%%%%%%%%%%%%%%%%%%%%%%%%%%%%%%%%%%%%%%%%%%%%%%%%%%%%%%%%%%%%%%%%%%%%%%%%%%%%%%%%%%%%
%%%%%%%%%%%%%%%%%%%%%%%%%%%%%%%%%%%%%%%%%%%%%%%%%%%%%%%%%%%%%%%%%%%%%%%%%%%%%%%%%%%%%%%%%%%%%%%%%%%%%%%%%%%

% \input{results}

\textbf{Declaration of generative AI and AI-assisted technologies in the writing process}

During the preparation of this work, the author(s) used ChatGPT (OpenAI) in order to assist with language editing, improving clarity, and refining the structure of the manuscript. After using this tool, the author(s) reviewed and edited the content as needed and take(s) full responsibility for the content of the publication.
\end{document}
